\chapter{Signals, Spectra and Systems}
\label{ch:signals}

\begin{chapterintro}
Strike a chord on a piano and two very different instruments will describe it. An oscilloscope
draws a jagged, restless line that never quite repeats. A spectrum analyzer draws three tidy
spikes, one for each note, each followed by a little family of overtones. Both pictures are
complete and both are correct; they are the same sound seen from two sides. A communication
engineer learns to flip between those two views without thinking, the way a bilingual person
switches languages mid-sentence, and that fluency is what this chapter is about.

It is the chapter that scares people most, so we will take it gently. We start with how big a
signal is and with the decibel, the engineer's Richter scale. We meet the single most useful idea
in the field---signals are arrows, with lengths and angles---and the spinning wheel that turns
every sinusoid into an arrow. Then Fourier's prism splits waveforms into pure tones, and we learn
what the transform's properties mean physically: delay is a phase slope, short is wide, filtering
is multiplication. We study the systems that carry signals---filters, cables, rooms full of
echoes---and what they do to a pulse, then the trick that turns a gigahertz radio signal into a
slow-moving complex number, the language the rest of the book is written in. We finish on the
practical side: what ``bandwidth'' really means, what an FFT plot actually shows, how a spectrum
analyzer sees the world, and what happens when an amplifier stops being linear. If you come from
computer engineering, read straight through; if you are an expert, the boxes hold the field lore
that textbooks usually leave out.
\end{chapterintro}

\mdcfig{ch02_two_domains}{The same sound in two languages. Left: a C-major chord (C, E and G,
each with a few overtones) as an oscilloscope shows it---a tangle with no obvious structure.
Right: the same 0.5\,s of signal as a spectrum analyzer shows it---three fundamentals and their
harmonics, every one sitting exactly where music theory says it should. Nothing was added or lost
between the two pictures; they are two coordinate systems for one signal.}

%======================================================================
\section{Signals as the language of communication}
\label{sec:ch02:signals}
%======================================================================

Every message you have ever sent has been a signal at some point: a wiggle of voltage in a
microphone cable, a swirl of current in a telegraph line, a shimmer of electric field at an
antenna, a flicker of light in a fibre, a stream of numbers pouring out of an analog-to-digital
converter. A \term{signal} is simply any quantity that varies with time (or space) and carries
information. Communication engineering is the art of shaping signals so that the information
survives a journey through a physical medium.

And the medium is never kind. It always does three things to a signal
(Figure~\ref{fig:ch02_three_insults}): it \emph{attenuates} it, it \emph{distorts} it (different
frequencies are delayed and weakened by different amounts), and it \emph{adds noise}. The tools of
this chapter describe the first two exactly; Chapter~\ref{ch:noise} deals with the third.

\mdcfig{ch02_three_insults}{The three insults every channel delivers. A clean bipolar bit pattern
(left) is first attenuated, then distorted by a channel that cannot keep up with fast
transitions---notice how the edges round off and the short pulses fail to reach full height---and
finally buried in noise. Everything in this book is a counter-attack on one of these three.}

\subsection{A taxonomy of signals}

Engineers sort signals along several independent axes, and each axis comes with its own
mathematics. Figure~\ref{fig:ch02_signal_zoo} is a field guide.

\mdcfig{ch02_signal_zoo}{A field guide to signals. Top: the same waveform in continuous time, as
discrete-time samples, and as digital samples rounded to a handful of levels. Bottom: a
deterministic cosine above a random noise record; a periodic square wave above a one-off pulse; and
a complex signal, which is best drawn not against time but as a path traced in the plane of its real
(I) and imaginary (Q) parts.}

\paragraph{Continuous time or discrete time.} A \emph{continuous-time} signal $x(t)$ is defined
for every real $t$: the voltage on a coaxial cable. A \emph{discrete-time} signal $x[n]$ is a
sequence of numbers, usually samples $x[n]=x(nT_s)$ taken every $T_s$ seconds. Every modern radio
is a continuous-time front end bolted to a discrete-time computer, and the boundary between
them---the sampling theorem of Chapter~\ref{ch:pcm}---is crossed in both directions millions of
times a second.

\paragraph{Analog or digital.} The \emph{values} of a signal may be continuous (analog) or drawn
from a finite set (digital). A sampled but unquantized signal is discrete-time and analog; the
output of an ADC is discrete-time and digital. A line-coded waveform on an Ethernet cable is
continuous-time but carries digital information.

\paragraph{Deterministic or random.} A \emph{deterministic} signal is a known function: a carrier
$\cos(2\pi f_c t)$, a pulse shape $p(t)$, a synchronisation preamble. A \emph{random} signal, like
noise or a stream of data symbols, can only be described statistically. Here is a lovely paradox:
information is inherently random---a signal you could predict would tell you nothing---so the
deterministic tools of this chapter are always applied, in the end, to the random processes of
Chapter~\ref{ch:noise}.

\paragraph{Periodic or aperiodic.} A periodic signal repeats, $x(t+T_0)=x(t)$ for all $t$; its
spectrum consists of discrete lines (Section~\ref{sec:ch02:fs}). An aperiodic signal has a
continuous spectrum (Section~\ref{sec:ch02:ft}).

\paragraph{Real or complex.} Every physical waveform is real. But two real waveforms---the
in-phase and quadrature components of a radio signal---are so naturally handled as one complex
signal that, from Section~\ref{sec:complexenvelope} onwards, almost every signal in this book is
complex. Do not let the word ``imaginary'' worry you; by the end of the chapter it will feel as
ordinary as a two-dimensional map coordinate.

\subsection{Energy and power}
\label{sec:ch02:energypower}

How \emph{big} is a signal? The question sounds naive, but it has two different answers depending
on whether the signal ends. A camera flash and a lighthouse both emit light; it makes sense to ask
how much energy the flash delivered, but the lighthouse has been burning for a century, so for it
the sensible question is how many watts it puts out on average.

Borrowing the convention of circuit theory, imagine $x(t)$ as a voltage across a 1\,$\Omega$
resistor. The instantaneous power is then $|x(t)|^2$, and the \term{energy} is
\begin{equation}
E_x=\int_{-\infty}^{\infty}|x(t)|^2\,dt .
\label{eq:ch02:energy}
\end{equation}
A signal with $0<E_x<\infty$ is an \term{energy signal}---the camera flash: a single pulse, a radar
chirp, a Wi-Fi packet viewed as a whole. Many signals of interest go on indefinitely and have
infinite energy; for these lighthouses the useful measure is the \term{average power}
\begin{equation}
P_x=\lim_{T\to\infty}\frac{1}{T}\int_{-T/2}^{T/2}|x(t)|^2\,dt ,
\label{eq:ch02:power}
\end{equation}
and a signal with $0<P_x<\infty$ is a \term{power signal}. Three answers are worth knowing by
heart. A sinusoid $A\cos(2\pi f_0 t+\phi)$ has power $A^2/2$, whatever its frequency and phase; a
constant (DC) level $A$ has power $A^2$; a complex exponential $Ae^{j2\pi f_0t}$ has power $|A|^2$,
because its magnitude never changes. A periodic signal's power is its energy over one period divided
by the period. Energy signals have zero average power and power signals have infinite energy, so
every signal of practical interest is one or the other.

\mdcpairany{ch02_peak_avg}{The radar worked example on a logarithmic power axis: brief 200\,W
spikes, an average of only 0.2\,W. The transmitter tube must survive the peaks; the electricity bill
sees only the average.}{ch02_scope}{An oscilloscope draws a signal the way the time-domain equations
of this chapter describe it: voltage against time. Here a triangle wave and a square wave share the
screen. By the end of Section~\ref{sec:ch02:fs} you will be able to say exactly which sinusoids each
is made of.}

The physical units follow from the 1\,$\Omega$ convention. If $x$ is a voltage across a resistance
$R$, divide by $R$; the shape of every calculation is unchanged. In digital signal processing the
``power'' of a sequence is simply its mean squared value, $P_x=\lim\frac{1}{N}\sum|x[n]|^2$, and
relating it to watts at an antenna requires a calibration chain that Chapter~\ref{ch:sdr}
describes.

\begin{worked}[Energy, power and the root-mean-square]
A radar transmits rectangular pulses of amplitude 100\,V across a 50\,$\Omega$ load, each
1\,$\mu$s long, at a pulse repetition frequency of 1\,kHz. Each pulse delivers energy
$E=(100^2/50)\times10^{-6}=200\,\mu$J. As a power signal, the pulse train has average power
$P=E\times\text{PRF}=200\,\mu\text{J}\times1000\,\text{s}^{-1}=0.2$\,W, while its \emph{peak}
power is $100^2/50=200$\,W (Figure~\ref{fig:ch02_peak_avg}). The ratio, $1000$ or 30\,dB, is the
reciprocal of the duty cycle. The root-mean-square (RMS) voltage is $\sqrt{0.2\times50}=3.2$\,V.
Peak-to-average power ratios reappear throughout the book: for radar they set transmitter design,
and for OFDM (Chapter~\ref{ch:ofdm}) they set power-amplifier efficiency.
\end{worked}

%======================================================================
\section{Decibels and engineering units}
\label{sec:ch02:decibels}
%======================================================================

In 1935 Charles Richter needed a way to compare earthquakes whose ground motions differed by
factors of millions, and he reached for a logarithm: each step on his scale is a factor of ten in
amplitude. Radio engineers had been doing the same thing for a decade. A cellular base station
transmits tens of watts; the phone a few kilometres away receives perhaps $10^{-12}$\,W; the
thermal noise in a 1\,Hz bandwidth is about $4\times10^{-21}$\,W. Multiplying and dividing numbers
with twenty-odd orders of magnitude is error-prone, so engineers work in logarithms. A power ratio
expressed in \term{decibels} is
\begin{equation}
G_{\mathrm{dB}}=10\log_{10}\frac{P_2}{P_1}.
\label{eq:ch02:db}
\end{equation}
In words: every factor of ten in power is 10\,dB. Because power is proportional to the square of
voltage (or field strength), a \emph{voltage} ratio in decibels is $20\log_{10}(V_2/V_1)$, provided
both voltages are measured across the same impedance. And here is the magic: gains in a cascade
multiply, so their decibel values \emph{add}. A link budget (Chapter~\ref{ch:noise} and
Chapter~\ref{ch:channels}) is nothing but a column of additions and subtractions.

\mdcpairany{ch02_seismogram}{A seismogram. Earthquakes span so many orders of magnitude that
seismologists measure them on a logarithmic scale; radio signals span even more, and engineers do
the same with the decibel.}{ch02_db_numbers}{The handful of decibel values that do almost all the
work, as power ratios (navy) and voltage ratios (red). Ten decibels is ten times the power but only
3.16 times the voltage.}

\mdcfig[0.8]{ch02_db_ruler}{The Richter scale for signals. From a 50\,kW broadcast
transmitter to the thermal noise in a single hertz of bandwidth is 251\,dB---a factor of more than
$10^{25}$---and every number on this ruler appears somewhere in this book.}

\begin{analogy}[The Richter scale for signals]
Nobody says ``that earthquake had 31\,623 times the ground motion of the last one''; they say ``it
was a 6.5, the last one was a 2''. Decibels do the same job for signals: they turn ``a millionth of
a millionth of a watt'' into $-90$\,dBm and turn multiplication into addition. Where the analogy
stops: the Richter scale is fixed to one reference, whereas a decibel is a pure \emph{ratio} until
you write a reference into its unit, as the next pages show.
\end{analogy}

\begin{history}[From ``miles of standard cable'' to the decibel]
Early telephone engineers measured loss in a delightfully concrete unit: the \emph{mile of
standard cable}, the attenuation of one mile of a particular reference telephone cable at voice
frequencies. It was awkward---it depended on frequency and on the cable---and in the 1920s the Bell
System replaced it with a logarithmic ``transmission unit'', soon renamed the \emph{decibel}: a
tenth of a \emph{bel}, a unit named in honour of Alexander Graham Bell. One decibel is roughly
the smallest change in loudness a listener can notice, which made it a natural step for telephone
work. A century later, the unit invented to compare telephone lines describes
everything from the output of a radio telescope to the noise inside a smartphone.
\end{history}

\subsection{A handful of numbers to know by heart}

A few decibel values do almost all the work. 3\,dB is a factor of 2 in power (more precisely
$10\log_{10}2=3.0103$); 10\,dB is a factor of 10; 1\,dB is about 26\%; 0.1\,dB is about 2.3\%. In
voltage, 6\,dB is a factor of 2 and 20\,dB a factor of 10. Everything else is built from these:
$17\,\text{dB}=10+10-3$ is a factor of $10\times10/2=50$, and $-13$\,dB is $1/20$. An engineer who
can do this in her head can check a link budget in a meeting---and will occasionally catch a
factor-of-a-thousand error before it is built into hardware.

\subsection{Absolute units: decibels relative to a reference}

A decibel is a \emph{ratio}. To express an absolute quantity, the reference is written into the
unit: dBm means ``decibels relative to one milliwatt'', so $+30$\,dBm is a watt and $-90$\,dBm is a
picowatt. Table~\ref{tab:ch02:dbunits} lists the units you will meet in this book and on instrument
screens; it is worth bookmarking.

\begin{table}[htbp]\centering\small
\caption{Decibel units in communications engineering.}\label{tab:ch02:dbunits}
\begin{tabular}{@{}lll@{}}
\toprule
Unit & Reference & Typical use\\
\midrule
dBW & 1 W & satellite EIRP, transmitter power \\
dBm & 1 mW & received power, instrument readings ($30\,\text{dBm}=1$\,W)\\
dBc & the carrier & spurs, harmonics, adjacent-channel leakage\\
dBc/Hz & carrier, per hertz & oscillator phase noise at an offset\\
dBm/Hz & 1 mW per hertz & noise density ($kT=-174$\,dBm/Hz at 290\,K)\\
dBFS & digital full scale & ADC and DSP signal levels\\
dB$\mu$V & 1 $\mu$V & receiver input voltage (0\,dBm in 50\,$\Omega$ = 107\,dB$\mu$V)\\
dB$\mu$V/m & 1 $\mu$V/m & field strength: EMC limits, broadcast coverage\\
dBi, dBd & isotropic, dipole & antenna gain ($0\,\text{dBd}=2.15\,\text{dBi}$)\\
dBK & 1 kelvin & system noise temperature in satellite $G/T$\\
dB-Hz & 1 Hz & carrier-to-noise density $C/N_0$ (GNSS, satellite)\\
\bottomrule
\end{tabular}
\end{table}

The density units deserve a moment, because they hide the most-used equation in radio. Thermal
noise power in a bandwidth $B$ is $kTB$, where $k=1.38\times10^{-23}$\,J/K is Boltzmann's constant
and $T$ the temperature. At the reference temperature of 290\,K, $kT=4.0\times10^{-21}$\,W/Hz,
which is $-204$\,dBW/Hz or $-174$\,dBm/Hz. A noise density in dBm/Hz becomes a power in dBm by
adding $10\log_{10}B$: in a 1\,MHz bandwidth the thermal noise floor is $-174+60=-114$\,dBm, and in a
20\,MHz Wi-Fi channel it is $-174+73=-101$\,dBm. This one line, with a receiver noise figure added,
is the most used equation in radio engineering (Chapter~\ref{ch:noise}).

\begin{worked}[Following a signal through a receiver]
A $-90$\,dBm signal enters an antenna port, passes a low-noise amplifier with 18\,dB gain, a
cable with 2\,dB loss, a mixer with 7\,dB conversion loss and an IF amplifier with 31\,dB
gain. The output level is $-90+18-2-7+31=-50$\,dBm, that is $10^{-5}$\,mW $=10$\,nW
(Figure~\ref{fig:ch02_db_cascade}). Across 50\,$\Omega$ this is an RMS voltage of
$\sqrt{10^{-8}\times50}=0.71$\,mV, or $-50+107=57$\,dB$\mu$V. A second $-90$\,dBm signal on another
frequency adds power, not decibels: the total is $10\log_{10}(2\times10^{-9}\,\text{mW})=-87$\,dBm,
not $-180$\,dBm.
\end{worked}

\mdcpairany{ch02_db_cascade}{A level diagram of the receiver worked example: each stage adds its
gain (or subtracts its loss) in decibels, and the level at any point is a running sum. RF engineers
draw exactly this picture on the whiteboard for every new design.}{ch02_specan}{A bench spectrum
analyzer. Its vertical axis is calibrated in dBm, its horizontal axis in frequency, and its
front-panel controls (span, resolution bandwidth, reference level) are all explained in
Section~\ref{sec:ch02:dft}.}

\begin{pitfall}[Ten or twenty?]
Use $10\log_{10}$ for ratios of powers, energies and power spectral densities, and
$20\log_{10}$ for ratios of amplitudes: voltages, currents, field strengths and the magnitude of a
transfer function $|H(f)|$. Plotting $20\log_{10}|X(f)|$ and $10\log_{10}|X(f)|^2$ gives the same
curve; plotting $20\log_{10}$ of a power spectral density doubles every number on the axis.
And never add two absolute levels in decibels: $-90\,\text{dBm}+(-90\,\text{dBm})$ is meaningless.
Add a level and a gain, or convert both levels to milliwatts first.
\end{pitfall}

%======================================================================
\section{Signals as vectors}
\label{sec:ch02:vectors}
%======================================================================

Here is the single most useful idea in communication theory, and it is worth meeting early:
\emph{signals are vectors}. Not as a metaphor, but in the precise sense that they can be added,
scaled, measured for length and compared by angle, with all the geometric intuition of arrows on a
sheet of paper. Once you see it, a whole shelf of apparently unrelated results---the matched filter,
OFDM, CDMA, MIMO, the Fourier series itself---collapses into one picture of a perpendicular dropped
from the tip of an arrow onto a line. Chapter~\ref{ch:modulation} turns this into the theory of
optimal detection; Fourier analysis, the subject of most of this chapter, is the special case in
which the arrows are sinusoids.

\mdcfig[0.55]{ch02_signal_distance}{Two signals as vectors. A sine and a square wave have lengths
(root energies) 0.71 and 0.85, are 0.37 apart, and have a correlation coefficient of 0.90---an angle
of $26^\circ$ between them. Every quantity in this caption is an inner product.}

\subsection{Inner products, length and angle}

For ordinary arrows, the dot product $\vect x\cdot\vect y=\sum_i x_iy_i^*$ multiplies matching
components and adds them up. A signal is like an arrow with infinitely many components---one for
every instant---so the sum becomes an integral. For complex-valued signals of finite energy define
the \term{inner product}
\begin{equation}
\langle x,y\rangle=\int_{-\infty}^{\infty}x(t)\,y^*(t)\,dt
\label{eq:ch02:inner}
\end{equation}
(over a finite interval when the signals live on one, as symbols do). It behaves exactly like a
dot product: it is linear in its first argument, it satisfies
$\langle y,x\rangle=\langle x,y\rangle^*$, and the inner product of a signal with itself is its
energy. So the \emph{length} or \emph{norm} of a signal is $\|x\|=\sqrt{\langle x,x\rangle}=\sqrt{E_x}$.
The distance between two signals is the norm of their difference,
$d(x,y)=\|x-y\|=\big(\int|x-y|^2dt\big)^{1/2}$---the root energy of the error. In a detector this
distance is the quantity that matters: the further apart two transmitted waveforms are, the more
noise it takes to confuse them.

The most important inequality in signal theory follows from these definitions. The
\term{Cauchy--Schwarz inequality} states that
\begin{equation}
|\langle x,y\rangle|\le\|x\|\,\|y\|,
\label{eq:ch02:cs}
\end{equation}
with equality if and only if $y$ is a scalar multiple of $x$. In words: two arrows can never
``agree'' more than the product of their lengths, and they reach that limit only when they point the
same way. It lets us define the angle between two signals through
$\cos\theta=\re\langle x,y\rangle/(\|x\|\|y\|)$, and its equality condition is the whole proof of
the matched filter (Chapter~\ref{ch:baseband}): the filter that maximises the output
signal-to-noise ratio is the one whose impulse response is ``parallel'' to the pulse it is looking
for. The normalised inner product $\rho=\langle x,y\rangle/(\|x\|\|y\|)$ is the
\term{correlation coefficient}; $|\rho|\le1$.

Two signals are \term{orthogonal} if $\langle x,y\rangle=0$---multiply them, add up the result,
and you get exactly nothing (Figure~\ref{fig:ch02_orthogonal_signals}). Orthogonal signals do not
interfere: the energy of their sum is the sum of their energies, $\|x+y\|^2=\|x\|^2+\|y\|^2$, which is
Pythagoras' theorem wearing a lab coat. Communication systems are built out of orthogonal signals:
pulses in different time slots (time-division multiplexing, Chapter~\ref{ch:multipleaccess});
sinusoids at frequencies that differ by a multiple of $1/T$ over a duration $T$ (the subcarriers of
OFDM, Chapter~\ref{ch:ofdm}); a cosine and a sine at the same frequency (the I and Q channels of
every modern modem, Section~\ref{sec:complexenvelope}); Walsh--Hadamard codes in CDMA
(Chapter~\ref{ch:spreadspectrum}); and the spatial signatures of users served by a multi-antenna base
station (Chapter~\ref{ch:mimo}). All five are the same idea in different coordinates.

\mdcfig{ch02_orthogonal_signals}{Four orthogonal pairs that run the modern world. In each column
the top two traces are the signals and the bottom trace is their product, shaded green where positive
and orange where negative. The green and orange areas cancel exactly, so the inner product is zero.
From left: the I and Q carriers of every radio; two TDMA time slots; two OFDM subcarriers; two CDMA
Walsh codes.}

\begin{analogy}[North and east]
Walk 3 km north and 4 km east and you are 5 km from home; no amount of walking north will ever
move you east. Orthogonal signals are north and east: a receiver that measures the ``northness'' of
what arrives is blind to anything sent ``east''. That is how a phone can listen to one OFDM
subcarrier while thousands of others shout in the same band. The analogy's limit: a signal space has
not two but infinitely many perpendicular directions---which is exactly why there is room for so many
users.
\end{analogy}

\subsection{Projection and best approximation}

Shine a torch straight down on a slanted stick and look at its shadow on the floor. The shadow is the
\emph{projection} of the stick onto the floor, and of all the arrows lying on the floor it is the one
closest to the stick. Signals work the same way. Suppose we want to approximate a signal $x$ by a
multiple $a\,y$ of another signal, choosing $a$ to minimise the energy of the error $\|x-ay\|^2$.
Expanding and setting the derivative to zero gives
\begin{equation}
a=\frac{\langle x,y\rangle}{\|y\|^2},
\label{eq:ch02:proj}
\end{equation}
the \term{projection} of $x$ on $y$. The error $x-ay$ is then orthogonal to $y$
(Figure~\ref{fig:ch02_projection}, left), just as the stick's ``height above its shadow'' is vertical.
If $\{\phi_k\}$ is a set of mutually orthogonal signals of unit energy---an \term{orthonormal} set,
$\langle\phi_k,\phi_m\rangle=\delta_{km}$---the best approximation of $x$ by a combination
$\sum_k a_k\phi_k$ uses the separately computed projections
\begin{equation}
a_k=\langle x,\phi_k\rangle ,
\label{eq:ch02:coeffs}
\end{equation}
and adding a new basis function never requires recomputing the old coefficients. This is the
property that makes orthogonal expansions so useful, and it is why the Fourier coefficients of
Section~\ref{sec:ch02:fs} can be computed one at a time. The right-hand panel of
Figure~\ref{fig:ch02_projection} projects a wiggly waveform onto the first three Legendre
polynomials, which are orthogonal on $[0,1]$; each added basis function lowers the error energy.

\mdcfig{ch02_projection}{Left: the projection of $\mathbf{x}$ on $\mathbf{y}$ is the multiple of
$\mathbf{y}$ closest to $\mathbf{x}$, and the error is perpendicular to $\mathbf{y}$. Right: the same
geometry with functions. The best approximations of $x(t)$ by one, two and three orthonormal
polynomials; the percentages are the fraction of the signal energy left in the error.}

When the orthonormal set is \emph{complete}---when the error can be made arbitrarily small for
every finite-energy signal---the expansion is exact in the energy sense, and the energy of the
signal equals the sum of the squared coefficients:
\begin{equation}
x(t)=\sum_k a_k\,\phi_k(t),\qquad E_x=\sum_k|a_k|^2 ,\qquad
\langle x,y\rangle=\sum_k a_k b_k^* .
\label{eq:ch02:parsevalgen}
\end{equation}
This is the general form of \term{Parseval's theorem}. Lengths and angles are preserved when we
replace a waveform by its list of coefficients; the waveform and the list are two coordinate systems
for the same vector, just as ``3 km north, 4 km east'' and ``5 km at bearing $53^\circ$'' describe one
journey. Section~\ref{sec:ch02:fs} uses sinusoids as the coordinates; Chapter~\ref{ch:modulation}
builds custom coordinates for a modulation scheme with the \term{Gram--Schmidt procedure}, which turns
any set of signals into an orthonormal basis by repeatedly subtracting projections.

\begin{keyidea}[Geometry is the shortest route to intuition]
Energy is squared length; correlation is angle; orthogonal signals do not interfere; the best
approximation is a projection and leaves a perpendicular error. Whenever a calculation in this
book becomes heavy---a matched filter, an equaliser, an OFDM receiver, a MIMO detector---ask what
the picture looks like in signal space. It almost always looks like
Figure~\ref{fig:ch02_projection}.
\end{keyidea}

%======================================================================
\section{Sinusoids, phasors and negative frequency}
\label{sec:ch02:phasors}
%======================================================================

Tie a reflector to the spoke of a bicycle wheel, spin the wheel, and watch it from the side with
your eye level with the axle. The reflector rises and falls, rises and falls, tracing out a perfect
sine wave in time (Figure~\ref{fig:ch02_bicycle_wheel}). That is the whole secret of this section:
\emph{a sinusoid is the shadow of something going round in a circle}. Keep the circle and throw away
the shadow, and every calculation with sinusoids becomes arithmetic with arrows.

\mdcfig{ch02_bicycle_wheel}{A phasor is a point on a spinning wheel. The arrow from the hub to the
reflector has a length (the amplitude $A$) and a starting angle (the phase $\phi$), and it turns at
$f_0$ revolutions per second. Its height, traced out over time, is a sinusoid. Engineers keep the
arrow and do arithmetic with it; the sinusoid can always be recovered as its shadow.}

The sinusoid $x(t)=A\cos(2\pi f_0 t+\phi)$ is the atom of signal analysis, for a simple reason:
linear time-invariant systems (Section~\ref{sec:lti}) map sinusoids to sinusoids of the same
frequency, changing only amplitude and phase. Three numbers describe it: the amplitude $A$, the
frequency $f_0$ in hertz (cycles per second), and the phase $\phi$ in radians. The period is
$T_0=1/f_0$ and the angular frequency is $\omega_0=2\pi f_0$; this book uses $f$ in hertz almost
everywhere because spectrum analyzers, regulators and datasheets do.

Euler's formula,
\begin{equation}
e^{j\theta}=\cos\theta+j\sin\theta ,
\end{equation}
is the bicycle wheel in mathematical form: as $\theta$ grows, $e^{j\theta}$ walks around the unit
circle, its real part is the horizontal shadow and its imaginary part the vertical one. It lets us
write the sinusoid as the real part of a rotating complex number, the \term{phasor}:
\begin{equation}
A\cos(2\pi f_0 t+\phi)=\re\{A e^{j\phi}e^{j2\pi f_0 t}\}
=\frac{A}{2}e^{j\phi}e^{j2\pi f_0 t}+\frac{A}{2}e^{-j\phi}e^{-j2\pi f_0 t}.
\label{eq:ch02:euler}
\end{equation}
The complex number $Ae^{j\phi}$ captures amplitude and phase in one quantity, and the factor
$e^{j2\pi f_0 t}$ simply turns it at $f_0$ revolutions per second. Adding sinusoids of the same
frequency reduces to adding complex numbers---adding arrows tip to tail
(Figure~\ref{fig:ch02_phasor_sum})---which is why AC circuit analysis, antenna-array theory and the
analysis of multipath fading (Chapter~\ref{ch:channels}) are all phasor arithmetic.

\mdcfig[0.5]{ch02_phasor_sum}{Two echoes of the same carrier arrive with different amplitudes and
phases. Their sum is found by placing the arrows tip to tail. Rotate the second arrow by
$180^\circ$---a path-length change of only half a wavelength---and the sum nearly vanishes: a
multipath fade.}

\mdcfig[0.85]{ch02_phasor}{Left: a real cosine is the sum of two phasors of half amplitude
rotating in opposite directions. Right: a phase offset $\phi$ rotates the phasor and shifts the
waveform in time. Carrier phase recovery (Chapter~\ref{ch:sync}) is the problem of estimating $\phi$.}

\subsection{Negative frequency: the wheel spinning backwards}

The second form of~\eqref{eq:ch02:euler}, shown in Figure~\ref{fig:ch02_phasor}, is the source of the
\emph{negative frequencies} that appear in every spectrum in this book, and which puzzle almost
everyone the first time. A real cosine is the sum of two counter-rotating phasors, one at $+f_0$ and
one at $-f_0$; their imaginary parts cancel and their real parts add.

So what \emph{is} a negative frequency? It is simply a wheel spinning backwards. Watch the wheel only
from the side and you cannot tell which way it turns: the shadow goes up and down identically either
way. That is a real signal, and it is why a real signal's spectrum is always symmetric. But look at
the wheel from the front---see both the horizontal and the vertical shadow at once---and clockwise
and counter-clockwise are obviously different. That is a complex signal, and a radio receiver with
I and Q outputs is exactly a pair of eyes looking at the wheel from the front. A complex signal
$e^{-j2\pi f_0 t}$ is a phasor rotating clockwise, and the receiver can tell it apart from one
rotating counter-clockwise: a tone 100\,kHz below the frequency to which it is tuned appears at
$-100$\,kHz, distinct from one 100\,kHz above (Figure~\ref{fig:ch02_neg_freq}).

\mdcfig{ch02_neg_freq}{Negative frequency made visible. Left: a complex (I and Q) signal containing
one tone at $+100$\,Hz and a weaker one at $-250$\,Hz; the two sides of the spectrum are different.
Right: keep only the real part and each tone acquires a mirror image of half its amplitude on the
other side of zero. A receiver with a single real ADC at zero IF cannot tell $+100$ from $-100$\,Hz;
one with I and Q can.}

For a \emph{real} signal, then, the spectrum at $-f$ is always the complex conjugate of the spectrum at
$+f$, so the negative-frequency half carries no new information. For a complex signal, both halves
are free. That freedom is the reason radios became IQ machines, and we shall cash it in
Section~\ref{sec:complexenvelope}.

\begin{analogy}[The wagon wheel in old westerns]
In old films the wagon wheels sometimes seem to turn slowly backwards while the wagon rolls forwards.
The camera samples the wheel 24 times a second, and a spoke that moves almost a full spoke-spacing
between frames looks as if it moved slightly backwards. The film camera sees only ``where the spoke
is now'', so it cannot tell a fast forward spin from a slow backward one. That is aliasing
(Chapter~\ref{ch:pcm}), and it is also a reminder that ``negative'' frequency is a perfectly physical
idea: a rotation in the other direction. The analogy's limit: in film the confusion is caused by
sampling; in a real-valued signal it happens even without sampling, because the side view alone
throws away the direction of spin.
\end{analogy}

\mdcfig[0.55]{ch02_wagon_wheel}{The wagon-wheel effect in numbers. A spoke turning 22 times a
second, filmed at 24 frames a second: the frames (dots) trace a slow two-revolution-per-second
motion, which the eye reads as the wheel creeping backwards.}

\begin{keyidea}[Think in phasors]
Multiplying a signal by $e^{j2\pi\Delta f t}$ shifts its whole spectrum by $\Delta f$; multiplying by
$e^{j\phi}$ rotates it by $\phi$. Almost every operation in a radio---mixing, tuning, carrier
recovery, Doppler correction, OFDM---is one of these two multiplications.
\end{keyidea}

\begin{tryit}[Spin the helix]
Open Lab~1 (\lab{lab01\_baseband.py}) and choose the \emph{IQ helix} experiment with the
\emph{Complex tone} signal. Drag the \emph{Frequency} slider from $+1$\,Hz through zero to
$-1$\,Hz and watch the helix change its sense of twist. Turn on the I and Q shadows: the shadows
look the same for both signs---only the pair together tells them apart. Then select
\emph{Real cosine} and watch it split into two counter-rotating helices.
\end{tryit}

%======================================================================
\section{Fourier series}
\label{sec:ch02:fs}
%======================================================================

In 1666 Isaac Newton darkened a room, let a thin beam of sunlight through a hole in the shutter, and
passed it through a glass prism. White light came out as a fan of colours. Newton's insight was that
the prism adds nothing: the colours were in the sunlight all along, and the glass merely sorts them
by bending each one by a different amount. A century and a half later Joseph Fourier claimed the
same thing for \emph{any} periodic waveform: it is a mixture of pure sinusoids, and a suitable
mathematical prism can sort them out. That prism is the Fourier series, and every spectrum analyzer
is a descendant of Newton's piece of glass.

\mdcphoto[0.4]{ch02_prism_photo}{Sunlight spread across a wall by a prism. The glass creates no
colours; it only sorts what was already in the white light by frequency. Fourier analysis does exactly
the same for signals.}

\mdcfig[0.95]{ch02_prism}{The Fourier prism. One awkward waveform (left) is exactly the sum of four
pure sinusoids at 1, 3, 5 and 7\,Hz (centre). The spectrum (right) records how much of each is
present: a bar chart of the ``colours'' in the signal.}

\subsection{Sinusoids as coordinates}

Now we can put the vector picture to work. Consider the complex exponentials $e^{j2\pi kf_0t}$,
$k=0,\pm1,\pm2,\dots$, on one period $T_0=1/f_0$. Any two of them are orthogonal over the period,
because their product is a whole number of cycles of a sinusoid, whose integral vanishes:
\begin{equation}
\frac{1}{T_0}\int_{T_0}e^{j2\pi kf_0t}\,e^{-j2\pi mf_0t}\,dt=
\begin{cases}1,&k=m,\\0,&k\ne m.\end{cases}
\end{equation}
So they form an orthogonal set---a set of perpendicular directions---and
by~\eqref{eq:ch02:coeffs} the best approximation of a periodic signal by a sum of them uses
coefficients equal to projections. The remarkable claim, made by Joseph Fourier in 1807, is that the
set is complete: \emph{every} periodic signal of finite power is a sum of harmonics of its
fundamental frequency. For a signal with period $T_0$,
\begin{equation}
x(t)=\sum_{k=-\infty}^{\infty}c_k\,e^{j2\pi k f_0 t},
\qquad c_k=\frac{1}{T_0}\int_{T_0}x(t)\,e^{-j2\pi k f_0 t}\,dt .
\label{eq:ch02:fs}
\end{equation}
Read the right-hand formula aloud: ``to find out how much of harmonic $k$ is in the signal, multiply
the signal by that harmonic (spinning backwards) and average over a period''. If the harmonic is
present, the product contains a steady component that survives the averaging; everything else
whirls around and averages to zero. The coefficient $c_k$ is the complex amplitude of the $k$-th
\term{harmonic}: $|c_k|$ is how much of it is present and $\angle c_k$ is where in the cycle it peaks.
The term $c_0$ is the average value (the DC component). For a real signal, $c_{-k}=c_k^*$, and
pairing the $\pm k$ terms gives the equivalent real form
\begin{equation}
x(t)=c_0+\sum_{k=1}^{\infty}2|c_k|\cos(2\pi kf_0t+\angle c_k).
\label{eq:ch02:fsreal}
\end{equation}
Plotting $|c_k|$ against $kf_0$ gives the \term{line spectrum} of the signal: a periodic signal
has power only at the harmonics of its fundamental, never between them.

\mdcpairany{ch02_graphic_eq}{A graphic equalizer: one slider per frequency band. Each slider scales
one ``coordinate'' of the music without touching the others, which only works because the bands are
(nearly) orthogonal.}{ch02_eq_bars}{The spectrum of a synthesised C-major chord (C3, C4, E4 and G4,
each with ten overtones) drawn the way an equalizer display draws music: power in half-octave bands.
The fundamentals fill the bands around 125--500\,Hz; the overtones climb towards 4\,kHz.}

\begin{analogy}[The bars on a music equalizer]
A stereo's graphic-equalizer display, with its columns of dancing lights, is a crude Fourier
analyzer: each column shows how much of the music lies in one band of frequencies at that moment
(Figure~\ref{fig:ch02_eq_bars}). Sliding a fader on the equalizer changes one coefficient and leaves
the others alone---exactly the independence that orthogonality guarantees. The analogy's limit: an
equalizer's bands are a fixed handful of octaves wide, while a Fourier series resolves each harmonic
exactly, and an equalizer display throws away the phase $\angle c_k$, which the series keeps.
\end{analogy}

\mdcphoto[0.36]{ch02_fourier}{Joseph Fourier, whose claim that any function could be built from
sines and cosines was greeted with scepticism by the greatest mathematicians of his day---and turned
out to underpin nearly everything in this book.}

\begin{history}[Fourier's heat, and the reception of an outrageous idea]
Joseph Fourier (1768--1830) was a mathematician, an administrator under Napoleon (he served as
prefect of the Is\`ere department) and a veteran of the Egyptian expedition, and his series came
out of a physics problem: how heat flows in a solid body. In December 1807 he presented a memoir
on heat propagation to the Institut de France, claiming that an arbitrary initial temperature
distribution---even a discontinuous one---could be expanded in sines and cosines. The examiners
included Lagrange, Laplace, Monge and Lacroix. Lagrange objected: he did not believe that a
discontinuous function could be a sum of smooth sinusoids, and the memoir was not published. Fourier
won the Institut's prize competition on heat in 1812 with a revised version, but the
judges still criticised its lack of rigour. His \emph{Th\'eorie analytique de la chaleur} finally
appeared in 1822. The doubts were justified in one sense---the convergence questions Fourier raised
took the rest of the century to settle, and they gave birth to much of modern analysis, from
Dirichlet's conditions to Riemann's integral and Cantor's set theory. But on the substance Fourier was
right, and the physicist's instinct that every signal is a sum of oscillations became the
foundation of acoustics, optics, quantum mechanics and every radio ever built.
\end{history}

\subsection{The square wave and the pulse train}

Two periodic waveforms appear so often that their series should be memorised. A square wave
that alternates between $+1$ and $-1$ with period $T_0$ has
\begin{equation}
x(t)=\frac{4}{\pi}\left[\sin(2\pi f_0t)+\frac13\sin(2\pi\,3f_0t)+\frac15\sin(2\pi\,5f_0t)+\cdots\right]:
\end{equation}
only odd harmonics, with amplitudes falling as $4/(\pi k)$ (Figure~\ref{fig:ch02_fourier_series}).
Watch it being built: one sinusoid is a rough sketch, three already show the shoulders, and fifty
are nearly indistinguishable from the square wave---except near the jumps.

\mdcfig{ch02_fourier_series}{Fourier synthesis of a square wave. The overshoot at the
discontinuity does not shrink as harmonics are added; it only narrows. The spectrum falls at
only 6\,dB per octave because the waveform has jumps.}

A train of rectangular pulses of height $A$, width $\tau$ and period $T_0$---the clock signal of every
digital circuit, a radar pulse train, a sampling waveform---has coefficients
\begin{equation}
c_k=A\,\frac{\tau}{T_0}\,\sinc\!\left(\frac{k\tau}{T_0}\right),\qquad
\sinc(x)=\frac{\sin\pi x}{\pi x}.
\label{eq:ch02:pulsetrain}
\end{equation}
The \term{duty cycle} $d=\tau/T_0$ sets the envelope: the lines sit at multiples of $f_0$ under a
$\sinc$ envelope whose first null is at $1/\tau$. With $d=\tfrac12$ the even harmonics fall on the
nulls of the envelope and vanish, which is why a perfect 50\% clock has no second harmonic;
any duty-cycle error brings the even harmonics back, and measuring the second harmonic of a clock
is a sensitive test of its symmetry.

\begin{tryit}[Build a square wave by hand]
In Lab~35 (\lab{lab35\_fourier\_spectra.py}), open the \emph{Fourier series builder}. With the
\emph{Square} target, drag \emph{Highest harmonic N} from 1 up to 99 and watch the overshoot at
the jump: it gets narrower but never lower. Switch on \emph{Lanczos $\sigma$-smoothing} and the
horns disappear---a preview of windowing. Then pick \emph{Pulse train} and slide the duty cycle to
exactly 0.5: every even harmonic drops out.
\end{tryit}

\subsection{Convergence and the Gibbs phenomenon}

In what sense does the series~\eqref{eq:ch02:fs} ``equal'' $x(t)$? For signals of finite power, the
partial sums $x_N(t)=\sum_{|k|\le N}c_ke^{j2\pi kf_0t}$ converge \emph{in energy}: the power of the
error $x-x_N$ tends to zero. That is the convergence an engineer usually needs, because it is the
error a receiver sees. Pointwise convergence needs more. The \term{Dirichlet conditions}---$x$
absolutely integrable over a period, with finitely many maxima, minima and discontinuities per
period---guarantee that the series converges to $x(t)$ wherever $x$ is continuous, and to the
midpoint of the jump at a discontinuity.

The convergence near a jump is not uniform, and the way it fails is famous. However many terms
are kept, the partial sums overshoot the jump by about 9\% of its height before settling
(Figure~\ref{fig:ch02_gibbs_zoom}). For the $\pm1$ square wave the peak of the partial sum
tends to $(2/\pi)\int_0^\pi(\sin u/u)\,du\approx1.179$, an overshoot of $0.179$ on a jump of 2. Adding
harmonics squeezes the overshoot closer to the discontinuity, so its energy goes to zero, but its
height does not. This is the \term{Gibbs phenomenon}.

\mdcpairany{ch02_gibbs_zoom}{A close-up of the jump. With 9, 49 and 199 harmonics the overshoot
peaks at 1.182, 1.179 and 1.179: it moves closer to the edge but never shrinks. Its energy goes to
zero; its height does not.}{ch02_gibbs}{Josiah Willard Gibbs of Yale, better known for founding
chemical thermodynamics, whose letters to \emph{Nature} explained the overshoot that now bears his
name.}

It is not a curiosity: it is what happens whenever a spectrum is truncated abruptly. A sharp
low-pass filter applied to a step rings with the same 9\% overshoot; the sidelobes of a rectangular
window (Section~\ref{sec:windows}) are the same effect seen from the other domain; and a video signal
passed through a brick-wall filter shows ``ringing'' ghosts beside every sharp edge. The cure is
always the same---taper the spectrum instead of truncating it---and the raised-cosine pulses of
Chapter~\ref{ch:baseband} and the windows of Section~\ref{sec:windows} are two applications of it.

\mdcpairany{ch02_michelson}{The Michelson--Stratton harmonic analyzer as illustrated in 1898. Each
of its many rocker arms contributes one sinusoid; springs add them, and a pen draws the
sum.}{ch02_michelson_gibbs}{Curves drawn by Michelson and Stratton's machine, reproduced from their
1898 paper: partial sums approaching a square wave. The small overshoots beside the sharp corners are
the Gibbs effect, drawn by gears and springs a century before anyone had a computer.}

\begin{history}[Michelson's machine and Gibbs's letter]
In 1898 Albert Michelson, already famous for measuring the speed of light, built with Samuel
Stratton a mechanical \emph{harmonic analyzer}: eighty springs and levers that could add eighty
Fourier components and draw the sum with a pen, or run in reverse to estimate the coefficients of a
drawn curve. When he fed it the coefficients of a square wave, the machine drew small ``horns''
beside each jump. Michelson suspected a defect of the instrument---or of Fourier's theory---and wrote
to \emph{Nature} questioning whether a sum of continuous functions could converge to a
discontinuous one. Josiah Willard Gibbs replied in \emph{Nature} in 1898 and, after correcting an
error, in 1899, explaining that the overshoot is a genuine property of the partial sums, not of the
machine. The effect had in fact been described by Henry Wilbraham in 1848, but his paper was
forgotten; Maxime B\^ocher named it after Gibbs in 1906. Michelson's analyzer survives, and its
drawings of the Gibbs horns---made with gears and springs before electronics existed---are among the
most charming artefacts of signal processing.
\end{history}

\subsection{Parseval's theorem and the power in harmonics}

Applying the general Parseval relation~\eqref{eq:ch02:parsevalgen} to the Fourier basis, the
power of a periodic signal is the sum of the powers of its harmonics:
\begin{equation}
P_x=\frac{1}{T_0}\int_{T_0}|x(t)|^2dt=\sum_{k=-\infty}^{\infty}|c_k|^2 .
\label{eq:ch02:parsevalfs}
\end{equation}
In words: you can add up the power either by squaring the waveform in time or by adding up the
power of each ``colour''---the books always balance. This is how distortion is quantified. The
\term{total harmonic distortion} (THD) of a nominally sinusoidal signal is the ratio of the RMS value
of all harmonics above the fundamental to the RMS value of the fundamental. Audio amplifiers are
specified with THD well below 0.1\%; a power amplifier driven hard into compression might show
several per cent.

\mdcfig[0.5]{ch02_thd_bars}{Where a square wave's power lives: 81\% in the fundamental, 9\% in the
third harmonic, 3\% in the fifth, and a long tail of ever smaller contributions.}

\begin{worked}[How much of a square wave is in its fundamental?]
A $\pm1$ square wave has power 1. Its fundamental has amplitude $4/\pi$ and therefore power
$\tfrac12(4/\pi)^2=8/\pi^2=0.811$. So 81\% of the power is in the fundamental and 19\% in the
harmonics (Figure~\ref{fig:ch02_thd_bars}). The THD is $\sqrt{(1-0.811)/0.811}=48\%$. The third
harmonic alone carries $0.811/9=9.0\%$ of the power and sits $20\log_{10}3=9.5$\,dB below the
fundamental, the fifth $14.0$\,dB below. These numbers explain why a saturated amplifier or a
switching mixer, which behaves like a square-wave multiplier, needs a harmonic filter after it, and
how much rejection that filter must provide.
\end{worked}

\subsection{Smoothness, decay and the spectrum of a clock}

How fast the harmonics decay is set by how smooth the waveform is
(Figure~\ref{fig:ch02_smoothness}). If $x$ has a jump, $|c_k|$ falls as $1/k$, or 20\,dB per decade of
frequency. If $x$ is continuous but its slope jumps (a triangle wave, or a trapezoid), $|c_k|$ falls
as $1/k^2$, 40\,dB per decade; each further derivative that is continuous buys another factor of $k$.
The reason is the differentiation property: differentiating $x$ multiplies $c_k$ by $j2\pi kf_0$, and
a function whose $n$-th derivative has jumps has coefficients that, multiplied by $k^n$, still fall
only as $1/k$. Smooth in time means compact in frequency. That one sentence explains why pulse
shaping (Chapter~\ref{ch:baseband}) works, why GSM uses Gaussian-filtered pulses
(Chapter~\ref{ch:modulation}), and why digital designers who care about emissions slow their clock
edges down.

\mdcfig{ch02_smoothness}{Smooth in time means compact in frequency. Left: a square wave (jumps), a
triangle wave (continuous, but with kinks) and a smooth periodic bump. Right: their harmonic
amplitudes on logarithmic axes. Jumps give a $1/k$ fall ($-20$\,dB per decade), kinks $1/k^2$
($-40$\,dB per decade), and the infinitely smooth bump falls off a cliff.}

A real clock is closer to a trapezoid than a square wave: it has finite rise and fall times $t_r$.
Its harmonic envelope has two corner frequencies. Below $1/(\pi\tau)$, where $\tau$ is the pulse
width, the envelope is flat; between $1/(\pi\tau)$ and $1/(\pi t_r)$ it falls at 20\,dB per
decade; above $1/(\pi t_r)$ it falls at 40\,dB per decade (Figure~\ref{fig:ch02_clock_emi}, left).
Slowing the edges from 200\,ps to 1\,ns moves the second corner from 1.6\,GHz down to 320\,MHz and
lowers the envelope above 1.6\,GHz by a factor of five, 14\,dB, without changing anything below
320\,MHz.

\begin{inpractice}[Clock harmonics and EMC compliance]
Every digital product sold in the United States must meet the radiated-emission limits of FCC
Part~15 (for consumer ``Class~B'' devices, 100\,$\mu$V/m at 3\,m from 30 to 88\,MHz, rising in
steps to 500\,$\mu$V/m above 960\,MHz); Europe and most other markets apply the similar CISPR~32
limits. Between 30\,MHz and 1\,GHz the measurement uses a receiver with a 120\,kHz resolution
bandwidth and a quasi-peak detector. The usual culprits are clock harmonics radiated by cables
and long traces, and the usual first-line fixes follow directly from Fourier analysis: slow the
edges (series termination resistors, controlled-slew drivers), keep duty cycles at 50\% to
suppress even harmonics, shorten the loops that radiate, and filter cables. The cleverest fix
is \term{spread-spectrum clocking} (SSC), now universal in PCs: the clock frequency is swept slowly,
for example by $-0.5$\% at a rate of 30--33\,kHz as PCI Express allows, so that each harmonic is
smeared over a band much wider than the 120\,kHz measurement bandwidth
(Figure~\ref{fig:ch02_clock_emi}, right). The total emitted power is unchanged, but the reading in
any one resolution bandwidth falls, typically by 5--15\,dB depending on the harmonic and the detector.
Higher harmonics are spread more, because a fractional deviation of 0.5\% is 0.5\% of a larger
frequency. The receivers at the other end of the link must tolerate the wandering clock, which is a
problem for the clock-recovery loops of Chapter~\ref{ch:sync}.
\end{inpractice}

\mdcfig{ch02_clock_emi}{Left: harmonics of a 100\,MHz, 50\%-duty clock with 1\,ns and 200\,ps
edges, with the asymptotic envelope (dashed) and the rise-time corner $1/(\pi t_r)$ (dotted).
Right: the ninth harmonic at 900\,MHz as an EMC receiver with a 120\,kHz resolution bandwidth
would see it, without and with 0.5\% down-spread spread-spectrum clocking (triangular profile at
31.5\,kHz). The energy is unchanged; it is spread across 4.5\,MHz, and the peak reading falls by
roughly 6\,dB with a peak detector and 14\,dB with an averaging detector.}

\begin{inpractice}[Harmonics are everywhere]
A transmitter's power amplifier driven into saturation turns a sinusoid into something
closer to a square wave, generating harmonics at $2f_c$, $3f_c$ and above. Regulators set
strict limits on these \emph{spurious emissions}---the FCC's general limit for many
licensed services is $43+10\log_{10}P$ decibels below the carrier, where $P$ is the transmit
power in watts---so every transmitter ends with a harmonic filter. The rate at which harmonics
decay is set by the smoothness of the waveform: a jump gives $1/k$ decay, a jump in slope
gives $1/k^2$, and so on. Section~\ref{sec:ch02:nonlinear} shows that a nonlinearity also produces
\emph{intermodulation} products close to the carrier, which no harmonic filter can remove.
\end{inpractice}

%======================================================================
\section{The Fourier transform}
\label{sec:ch02:ft}
%======================================================================

A drummer hits a snare once. There is no period, nothing repeats, so the Fourier series has nothing
to say. Yet the snare clearly has a ``sound''---a spread of frequencies we can hear. The Fourier
transform is the Fourier series grown up: the prism that works on signals that happen only once.

\mdcfig[0.85]{ch02_ft_limit}{From Fourier series to Fourier transform. The pulse train with pulse
width $\tau$ and period $T_0$ has lines every $1/T_0$; scaled by $T_0$, they sample the fixed
envelope $\tau\sinc(f\tau)$ (dashed), which is the Fourier transform of a single pulse. As
$T_0\to\infty$ the samples fill in the curve.}

\subsection{From series to transform}

The trick is to pretend. A single pulse can be thought of as one period of a periodic signal whose
period is enormous---so long that the next pulse never arrives. Figure~\ref{fig:ch02_ft_limit}
shows what happens to the line spectrum of a pulse train as the period $T_0$ grows with the pulse
width $\tau$ fixed. The lines, spaced $f_0=1/T_0$ apart, crowd together, and each line gets weaker
(the power is shared among more of them), but when scaled by $T_0$ they trace out a fixed envelope.
In the limit, the lines merge into a continuous curve and we obtain the \term{Fourier transform}
pair
\begin{equation}
X(f)=\int_{-\infty}^{\infty}x(t)\,e^{-j2\pi f t}\,dt,
\qquad x(t)=\int_{-\infty}^{\infty}X(f)\,e^{j2\pi f t}\,df.
\label{eq:ch02:ft}
\end{equation}

\begin{deeper}[The limit, step by step]
Write the Fourier series~\eqref{eq:ch02:fs} as
\begin{equation*}
T_0c_k=\int_{-T_0/2}^{T_0/2}x(t)\,e^{-j2\pi kf_0t}\,dt,\qquad
x(t)=\sum_k (T_0c_k)\,e^{j2\pi kf_0t}\,f_0 .
\end{equation*}
Now let $T_0\to\infty$. The harmonic frequency $kf_0$ becomes a continuous frequency $f$, the line
spacing $f_0$ becomes $df$, the quantity $T_0c_k$ becomes the function $X(f)$, and the sum over $k$
becomes an integral over $f$. The two equations turn into~\eqref{eq:ch02:ft}.
\end{deeper}

We write $x(t)\leftrightarrow X(f)$. The first integral is the \emph{analysis} equation: it projects
$x$ onto the complex exponential at frequency $f$, exactly as~\eqref{eq:ch02:coeffs} did for a
discrete basis---``how much of frequency $f$ is in this signal?''. The second is the \emph{synthesis}
equation: it rebuilds $x$ as a continuous superposition of exponentials. A periodic signal has a
spectrum of discrete lines with finite power in each; an aperiodic signal has a continuous spectrum,
and $X(f)$ is a \emph{density}, with units of (signal units)$\times$seconds, or equivalently per hertz.

We use ordinary frequency $f$ in hertz rather than angular frequency $\omega$ because it
removes all the stray factors of $2\pi$ from the transform pair and because engineers read
spectra in hertz. $X(f)$ is complex: its magnitude says how much of each frequency is present,
and its phase says how the components are aligned in time. Two signals with the same
$|X(f)|$ but different phases can look utterly different---an impulse and a chirp have the same
flat magnitude spectrum---so a magnitude plot alone never tells the whole story.

The transform exists in the ordinary sense for absolutely integrable signals, and in an energy
sense (the Plancherel theorem) for every finite-energy signal. Power signals such as sinusoids
and constants have no transform in the ordinary sense, but they acquire one when we admit the
impulse $\delta(f)$, as Section~\ref{sec:ch02:delta} shows.

\subsection{Properties worth knowing by heart}

Table~\ref{tab:ch02:props} lists the properties used most often. Do not memorise it as algebra;
each line is a piece of physics, and the paragraphs after it say what each one \emph{means}.

\begin{table}[htbp]\centering\small
\caption{Fourier transform properties ($x(t)\leftrightarrow X(f)$).}\label{tab:ch02:props}
\begin{tabular}{@{}lll@{}}
\toprule
Property & Time domain & Frequency domain\\
\midrule
Linearity & $a\,x(t)+b\,y(t)$ & $a\,X(f)+b\,Y(f)$\\
Time shift & $x(t-\tau)$ & $e^{-j2\pi f\tau}X(f)$\\
Frequency shift & $x(t)\,e^{j2\pi f_0 t}$ & $X(f-f_0)$\\
Modulation & $x(t)\cos(2\pi f_0 t)$ & $\tfrac12[X(f-f_0)+X(f+f_0)]$\\
Scaling & $x(at)$ & $\frac{1}{|a|}X(f/a)$\\
Time reversal & $x(-t)$ & $X(-f)$\\
Conjugation & $x^*(t)$ & $X^*(-f)$\\
Convolution & $x(t)*y(t)$ & $X(f)\,Y(f)$\\
Multiplication & $x(t)\,y(t)$ & $X(f)*Y(f)$\\
Differentiation & $\frac{d}{dt}x(t)$ & $j2\pi f\,X(f)$\\
Integration & $\int_{-\infty}^{t}x(s)\,ds$ & $\frac{X(f)}{j2\pi f}+\frac12X(0)\delta(f)$\\
Duality & $X(t)$ & $x(-f)$\\
Area & $\int x(t)\,dt=X(0)$ & $x(0)=\int X(f)\,df$\\
Parseval (Rayleigh) & $\int|x(t)|^2dt$ & $\int|X(f)|^2df$\\
\bottomrule
\end{tabular}
\end{table}

\paragraph{Time shift: delay is phase slope.} Delaying a signal by $\tau$ leaves $|X(f)|$ alone
and adds a phase $-2\pi f\tau$ that grows linearly with frequency (Figure~\ref{fig:ch02_delay_phase}).
Why? A delay of $\tau$ is a small fraction of a cycle at low frequency and many cycles at high
frequency. A delay of 1\,ns is a phase shift of $360^\circ$ at 1\,GHz and of $3.6^\circ$ at 10\,MHz.
Conversely, a phase that is linear in frequency is nothing but a delay, and a phase that is
\emph{not} linear is distortion (Section~\ref{sec:ch02:distortionless}).

\paragraph{Frequency shift: multiplying by a phasor moves the spectrum.} Multiplying by
$e^{j2\pi f_0t}$ slides the whole spectrum up by $f_0$. This is the mathematics of every
transmitter, every tuner and every digital down-converter, and it is the frequency-domain twin of
the time-shift property.

\paragraph{Scaling: short pulses need wide bandwidth.} Compressing a signal in time by a factor
$a$ stretches its spectrum by $a$ and lowers it by $a$. Doubling a data rate doubles the
bandwidth; a 1\,ns pulse occupies roughly a gigahertz. A snare hit, over in a few milliseconds,
sprays energy across a wide band; a sustained violin note is a narrow line
(Figure~\ref{fig:ch02_drum_violin}). This reciprocity, sharpened into an inequality in
Section~\ref{sec:ch02:uncertainty}, is the most important single fact about spectra.

\mdcfig[0.5]{ch02_delay_phase}{A pure delay leaves the magnitude alone and adds a phase that falls in a
straight line with frequency. The slope is the delay: twice the delay, twice the slope.}

\paragraph{Convolution becomes multiplication.} The output of a linear time-invariant system is
the convolution of the input with the impulse response (Section~\ref{sec:lti}); in frequency, it
is a product. Filters are therefore designed and understood in the frequency domain, and fast
convolution with the FFT (Chapter~\ref{ch:dsp}) and the cyclic prefix of OFDM
(Chapter~\ref{ch:ofdm}) are both applications of this property.

\mdcfig{ch02_drum_violin}{Short in time, wide in frequency. A drum-like hit at 440\,Hz that dies away
in a few milliseconds (top) has a spectrum 79\,Hz wide at the 3\,dB points; a violin-like note at the
same pitch held for 0.8\,s (bottom) is only 2\,Hz wide. Same pitch, a forty-fold difference in
bandwidth, set entirely by duration.}

\paragraph{Multiplication becomes convolution.} Multiplying two signals convolves their spectra,
so the bandwidth of a product is roughly the sum of the bandwidths. Windowing a signal before an
FFT (Section~\ref{sec:windows}), gating a carrier into a pulse, and sampling (multiplying by an
impulse train) all smear or replicate the spectrum in this way.

\paragraph{Differentiation emphasises high frequencies.} Each derivative multiplies the spectrum
by $j2\pi f$, a 20\,dB-per-decade boost. That is why differentiating a noisy signal is a bad idea,
why FM discriminators produce noise that rises with frequency (Chapter~\ref{ch:analog}), and why
smooth signals have rapidly decaying spectra.

\paragraph{Duality.} The two integrals in~\eqref{eq:ch02:ft} differ only in the sign of the
exponent, so every transform pair yields a second pair with the roles of time and frequency
exchanged. A rectangular pulse has a sinc spectrum, so a sinc pulse has a rectangular spectrum. An
impulse has a flat spectrum, so a constant has an impulsive spectrum. Half of the table of pairs
comes free.

\paragraph{Parseval's theorem.} Energy can be computed in either domain,
\begin{equation}
E_x=\int_{-\infty}^{\infty}|x(t)|^2dt=\int_{-\infty}^{\infty}|X(f)|^2df ,
\label{eq:ch02:parseval}
\end{equation}
and more generally $\langle x,y\rangle=\langle X,Y\rangle$: the Fourier transform preserves lengths
and angles. It is a rotation of signal space, the infinite-dimensional analogue of choosing new
coordinate axes. Orthogonal signals stay orthogonal, and an error of a given energy in one domain
is an error of the same energy in the other.

\paragraph{Symmetry of real signals.} If $x$ is real, the conjugation property gives
$X(-f)=X^*(f)$: the magnitude spectrum is even and the phase is odd. A real signal's spectrum is
determined by its positive-frequency half---the bicycle wheel seen from the side---a fact that
Section~\ref{sec:complexenvelope} exploits.

\subsection{Transform pairs}

A few pairs appear over and over (Figure~\ref{fig:ch02_pairs} and Table~\ref{tab:ch02:pairs}), and
together they are the vocabulary of the rest of the book. The rectangular pulse of width $T$
transforms to a sinc function,
\begin{equation}
\operatorname{rect}(t/T)\ \leftrightarrow\ T\,\sinc(fT),
\label{eq:ch02:rectsinc}
\end{equation}
whose first nulls are at $\pm1/T$ and whose sidelobes decay slowly, as $1/f$, because the pulse has
jumps. By duality, an ideal ``brick-wall'' low-pass spectrum corresponds to a sinc pulse in
time---the pulse Nyquist proposed for signalling without intersymbol interference
(Chapter~\ref{ch:baseband}). The triangle, the convolution of two rectangles, transforms to
$\sinc^2$, with sidelobes decaying as $1/f^2$ because the triangle is continuous with jumps only in
slope. The Gaussian is its own transform, and among all pulses it minimises the product of
duration and bandwidth (Section~\ref{sec:ch02:uncertainty}). The one-sided exponential, the impulse
response of an $RC$ circuit, transforms to a Lorentzian whose magnitude falls as $1/f$.

\mdcfig[0.85]{ch02_pairs}{Four Fourier transform pairs. Smoother time functions have spectra that
decay faster; narrower time functions have wider spectra.}

\begin{table}[htbp]\centering\small
\caption{Common Fourier transform pairs. $\operatorname{rect}(t)$ is 1 for $|t|<\tfrac12$ and 0
otherwise; $\Lambda(t)=\max(1-|t|,0)$ is the unit triangle; $u(t)$ is the unit step.}
\label{tab:ch02:pairs}
\begin{tabular}{@{}lll@{}}
\toprule
Signal & $x(t)$ & $X(f)$\\
\midrule
Rectangular pulse & $\operatorname{rect}(t/T)$ & $T\sinc(fT)$\\
Sinc pulse & $2W\sinc(2Wt)$ & $\operatorname{rect}(f/2W)$\\
Triangle & $\Lambda(t/T)$ & $T\sinc^2(fT)$\\
Gaussian & $e^{-\pi t^2/T^2}$ & $T\,e^{-\pi f^2T^2}$\\
One-sided exponential & $e^{-at}u(t)$, $a>0$ & $1/(a+j2\pi f)$\\
Two-sided exponential & $e^{-a|t|}$ & $2a/(a^2+4\pi^2f^2)$\\
Impulse & $\delta(t)$ & 1\\
Constant & 1 & $\delta(f)$\\
Complex exponential & $e^{j2\pi f_0t}$ & $\delta(f-f_0)$\\
Cosine & $\cos(2\pi f_0t)$ & $\tfrac12[\delta(f-f_0)+\delta(f+f_0)]$\\
Sine & $\sin(2\pi f_0t)$ & $\tfrac1{2j}[\delta(f-f_0)-\delta(f+f_0)]$\\
Signum & $\operatorname{sgn}(t)$ & $1/(j\pi f)$\\
Unit step & $u(t)$ & $\tfrac12\delta(f)+1/(j2\pi f)$\\
Impulse train & $\sum_n\delta(t-nT)$ & $\frac1T\sum_k\delta(f-k/T)$\\
\bottomrule
\end{tabular}
\end{table}

\subsection{Modulation and frequency translation}
\label{sec:ch02:modthm}

How does a voice get onto a radio carrier at 100\,MHz, and how does it get off again? By
multiplication, and the Fourier transform tells us exactly what multiplication does. The
\term{modulation theorem} is the frequency-shift property applied to a real carrier:
\begin{equation}
x(t)\cos(2\pi f_ct)\ \leftrightarrow\ \tfrac12X(f-f_c)+\tfrac12X(f+f_c).
\label{eq:ch02:modthm}
\end{equation}
Multiplying a baseband signal by a carrier makes two half-size copies of its spectrum, one centred
on $+f_c$ and one on $-f_c$ (Figure~\ref{fig:ch02_mod_shift}). This is double-sideband modulation
(Chapter~\ref{ch:analog}), and it is also what every mixer does. Multiplying again by
$\cos(2\pi f_ct)$ moves each copy up and down by $f_c$ once more, giving a copy at zero frequency and
copies at $\pm2f_c$; a low-pass filter keeps the first, and the original signal is recovered. That is
coherent demodulation. If the receiver's carrier has a phase error $\phi$, the recovered signal is
scaled by $\cos\phi$---a $90^\circ$ error makes it vanish---which is why carrier recovery matters
(Chapter~\ref{ch:sync}).

\mdcpairany{ch02_mod_shift}{The modulation theorem in three steps. Multiplying by a carrier makes
two half-size copies at $\pm f_c$; multiplying again brings one copy home to zero frequency and
throws others out to $\pm2f_c$; a low-pass filter keeps the one at home.}{ch02_sampling_comb}{Sampling
is multiplication by an impulse train, so it replicates the spectrum every $f_s$. If the copies do
not overlap, nothing is lost---the sampling theorem of Chapter~\ref{ch:pcm} in one picture.}

A mixer multiplies two signals, and the product of two sinusoids is a sum and a difference:
\begin{equation}
\cos(2\pi f_1t)\cos(2\pi f_2t)=\tfrac12\cos\big(2\pi(f_1-f_2)t\big)+\tfrac12\cos\big(2\pi(f_1+f_2)t\big).
\end{equation}
A superheterodyne receiver with its local oscillator at $f_{\mathrm{LO}}$ converts both
$f_{\mathrm{LO}}+f_{\mathrm{IF}}$ and $f_{\mathrm{LO}}-f_{\mathrm{IF}}$ to the same intermediate frequency
$f_{\mathrm{IF}}$; the unwanted one is the \term{image}, and rejecting it is the central problem of
receiver design (Chapter~\ref{ch:analog} and Chapter~\ref{ch:sdr}). A complex mixer, multiplying by
$e^{-j2\pi f_{\mathrm{LO}}t}$, shifts in only one direction and has no image in principle; that is
the real reason radios became IQ machines (Section~\ref{sec:complexenvelope}).

\begin{tryit}[Lose the signal on purpose]
In Lab~1's \emph{Up- and down-conversion} experiment, choose the two-tone message and drag the
\emph{LO phase error} slider. The recovered tones shrink as $\cos\phi$ and vanish at $\pm90^\circ$.
Now set a small \emph{LO frequency error} instead: the output swells and fades periodically---the
phase error is now growing with time.
\end{tryit}

\begin{worked}[The spectrum of a radar pulse]
A radar transmits a 1\,$\mu$s rectangular burst of a 3\,GHz carrier: $x(t)=\operatorname{rect}(t/T)\cos(2\pi f_ct)$
with $T=1\,\mu$s. By~\eqref{eq:ch02:rectsinc} and~\eqref{eq:ch02:modthm},
$X(f)=\tfrac T2\sinc\big((f-f_c)T\big)+\tfrac T2\sinc\big((f+f_c)T\big)$. Around 3\,GHz the spectrum is a
sinc with nulls at $3000\pm1$\,MHz, a main lobe 2\,MHz wide, and first sidelobes 13.3\,dB down. The
receiver needs a bandwidth of about $1/T=1$\,MHz to pass most of the energy. The range resolution
is $c\,T/2=150$\,m: two targets closer than that return overlapping echoes. A finer resolution
seems to demand a shorter pulse and therefore more peak power for the same energy---but
Section~\ref{sec:ch02:correlation} shows how pulse compression escapes the trade by making the
bandwidth large without making the pulse short.
\end{worked}

\mdcfig[0.55]{ch02_radar_spectrum}{The radar pulse of the worked example around its 3\,GHz carrier:
a sinc-shaped spectrum with nulls 1\,MHz either side of the carrier and first sidelobes 13.3\,dB
down.}

\subsection{Impulses and the transforms of power signals}
\label{sec:ch02:delta}

Strike a tuning fork with a rubber mallet and it rings at its one natural frequency; strike a bell
and you hear all its modes at once. A sharp tap is the universal test signal, because it contains
every frequency. The mathematical idealisation of that tap is the Dirac \term{impulse} (delta
function) $\delta(t)$: a pulse so short and so tall that only its area, 1, matters. It is defined by
what it does inside an integral, the \emph{sifting property}
\begin{equation}
\int_{-\infty}^{\infty}x(t)\,\delta(t-t_0)\,dt=x(t_0),
\end{equation}
and it can be pictured as the limit of rectangles $\frac1\epsilon\operatorname{rect}(t/\epsilon)$
or Gaussians of unit area as their width goes to zero. Mathematicians call it a
\emph{generalized function} or \emph{distribution}; for engineers, it is enough to know that every
step involving it can be justified by working with a narrow pulse and taking the limit at the end.

\mdcfig[0.75]{ch02_impulse_flat}{Squeeze a pulse while keeping its area fixed and its spectrum
flattens. In the limit---the impulse---every frequency is present equally, which is why a sharp tap
reveals a system's whole frequency response.}

\mdcphoto[0.4]{ch02_tuning_forks}{Tuning forks: tap one and it rings at the single frequency its
shape selects. The tap contains every frequency; the fork's response picks one. The larger the fork,
the lower its note.}

Sifting gives $\delta(t)\leftrightarrow1$: an impulse contains all frequencies equally, with zero
phase, which is why striking a bell (or a channel) with an impulse reveals its whole frequency
response. Duality gives $1\leftrightarrow\delta(f)$: a constant is pure DC. The frequency-shift
property then gives the transforms of all sinusoids, $e^{j2\pi f_0t}\leftrightarrow\delta(f-f_0)$
and $\cos(2\pi f_0t)\leftrightarrow\tfrac12[\delta(f-f_0)+\delta(f+f_0)]$: a spectral line is an
impulse in frequency. A periodic signal is a sum of exponentials~\eqref{eq:ch02:fs}, so its
transform is a train of impulses weighted by its Fourier coefficients,
\begin{equation}
X(f)=\sum_{k=-\infty}^{\infty}c_k\,\delta(f-kf_0),
\label{eq:ch02:periodicft}
\end{equation}
and the Fourier series turns out to be a special case of the Fourier transform. One periodic signal
deserves special mention: the impulse train $\sum_n\delta(t-nT)$ transforms to an impulse train
$\frac1T\sum_k\delta(f-k/T)$, spaced $1/T$ apart. Multiplying a signal by an impulse train is
sampling, and by the multiplication property it replicates the spectrum every $1/T$\,Hz
(Figure~\ref{fig:ch02_sampling_comb}). That single observation is the sampling theorem and the theory
of aliasing, which Chapter~\ref{ch:pcm} develops in full; Section~\ref{sec:ch02:dft} uses it to
explain the DFT.

\subsection{Duration, bandwidth and the uncertainty principle}
\label{sec:ch02:uncertainty}

Musicians know this one in their bones. A drum hit is over in an instant, and you could not hum its
pitch; a long bowed violin note has a pitch you could tune a piano to. The shorter a sound, the
less definite its frequency. That is not a limitation of ears or instruments---it is a law of
mathematics, and in communications it sets the price of speed: every pulse you make shorter, to send
data faster, costs you bandwidth.

\mdcfig{ch02_uncertainty}{The uncertainty principle. Left: four pulses of unit energy and similar
RMS duration, with their time--bandwidth products. Right: their spectra. The Gaussian meets the
bound $\sigma_t\sigma_f=1/(4\pi)$ and has the fastest-decaying spectrum; the rectangle's spectrum
decays so slowly that its RMS bandwidth is infinite.}

\begin{analogy}[A drum hit and a violin note]
Ask ``what frequency is this sound?'' and the honest answer depends on how long you are allowed to
listen. To pin down a frequency to within 1\,Hz you must hear about a second of it, because you need
to count cycles; a sound that lasts a millisecond simply does not contain enough cycles to be precise
about. The drum hit of Figure~\ref{fig:ch02_drum_violin} is not ``out of tune''; it genuinely has no
sharp frequency. The analogy's limit: musicians talk about pitch, a perceptual quantity; the
uncertainty principle is about the spread of energy, which is exactly measurable.
\end{analogy}

The \term{uncertainty principle} makes this exact. Treat $|x(t)|^2/E_x$ and $|X(f)|^2/E_x$ as
probability densities (they are non-negative and, by Parseval, both integrate to 1) and measure the
spread of each by its standard deviation, the \term{RMS duration} $\sigma_t$ and the
\term{RMS bandwidth} $\sigma_f$:
\begin{equation}
\sigma_t^2=\frac1{E_x}\int(t-\bar t)^2|x(t)|^2dt,\qquad
\sigma_f^2=\frac1{E_x}\int(f-\bar f)^2|X(f)|^2df .
\end{equation}
Then for every signal
\begin{equation}
\sigma_t\,\sigma_f\ \ge\ \frac{1}{4\pi},
\label{eq:ch02:uncertainty}
\end{equation}
with equality only for a Gaussian pulse (possibly shifted in time and frequency). In words: you can
squeeze a pulse in time or in frequency, but the product of the two spreads can never fall below
about 0.08, and only the bell-shaped Gaussian reaches that floor.

\begin{deeper}[Proof of the uncertainty principle]
The proof is a neat use of the tools above. Take $\bar t=\bar f=0$. By Parseval and the
differentiation property, $\int|x'(t)|^2dt=4\pi^2\int f^2|X(f)|^2df=4\pi^2E_x\sigma_f^2$. By the
Cauchy--Schwarz inequality~\eqref{eq:ch02:cs} applied to $t\,x(t)$ and $x'(t)$,
\[
\Big|\int t\,x\,x'^*\,dt\Big|^2\le\int t^2|x|^2dt\cdot\int|x'|^2dt=E_x\sigma_t^2\cdot4\pi^2E_x\sigma_f^2 .
\]
Integrating by parts shows that the real part of the left-hand integral is $-E_x/2$, so
$E_x^2/4\le4\pi^2E_x^2\sigma_t^2\sigma_f^2$, which is~\eqref{eq:ch02:uncertainty}. Equality in
Cauchy--Schwarz requires $x'(t)\propto t\,x(t)$, whose solution is the Gaussian. The same mathematics
is Heisenberg's uncertainty relation $\sigma_x\sigma_p\ge h/(4\pi)$ between position and momentum,
because a particle's momentum is Planck's constant $h$ times the spatial frequency of its wavefunction.
\end{deeper}

Figure~\ref{fig:ch02_uncertainty} compares four unit-energy pulses of similar RMS duration. The
Gaussian achieves $\sigma_t\sigma_f=1/(4\pi)\approx0.080$; a raised-cosine (Hann-shaped) pulse is
within 3\% of the bound at 0.082; a half-sine pulse, the shape used in MSK
(Chapter~\ref{ch:modulation}), reaches 0.090. The rectangle has no finite RMS bandwidth at all: its
$\sinc^2$ power spectrum decays as $1/f^2$, so $\int f^2|X|^2df$ diverges. That is not a quibble. It
is a warning that ``bandwidth'' has no meaning until it is defined
(Section~\ref{sec:bandwidth}), and that a pulse with jumps spreads energy over the whole spectrum.

\begin{tryit}[Squeeze a pulse]
In Lab~35, open \emph{The uncertainty principle}. Pick the \emph{Gaussian} and drag \emph{Pulse
width} from 2\,s down to 0.25\,s: the spectrum widens in exact proportion and the dot on the
$\sigma_t$--$\sigma_f$ map slides along the bound. Switch to \emph{Raised cosine}, \emph{Half sine}
and \emph{Rectangle} and see how far above the bound each one sits.
\end{tryit}

\mdcphoto[0.34]{ch02_gabor}{Dennis Gabor in 1971, the year of his Nobel Prize. His 1946 paper gave
communication engineers the uncertainty principle, the time--frequency plane and the analytic
signal---three ideas this chapter cannot do without.}

\begin{inpractice}[Gaussian pulses in GSM, Bluetooth and DECT]
Because the Gaussian is the most compact pulse in both domains at once, it is the natural choice
when spectrum must be kept tight around a constant-envelope signal. GSM transmits with
\emph{Gaussian minimum-shift keying} (GMSK): the data are filtered by a Gaussian filter whose
3\,dB bandwidth $B$ is $0.3$ times the bit rate (the standard's $BT=0.3$) before driving a frequency
modulator. Bluetooth Basic Rate uses Gaussian frequency-shift keying with $BT=0.5$, and DECT
cordless phones use $BT=0.5$ as well. A smaller $BT$ gives a narrower spectrum but stretches each
pulse over more bit periods, creating intersymbol interference that the receiver must handle. The
trade-off is exactly~\eqref{eq:ch02:uncertainty}, set to a particular value by the system
designer (Chapter~\ref{ch:modulation}).
\end{inpractice}

\begin{history}[Gabor, the logon and the analytic signal]
In 1946 Dennis Gabor, a Hungarian-born engineer working for British Thomson-Houston in Rugby,
published ``Theory of Communication'' in the \emph{Journal of the Institution of Electrical
Engineers}. Physicists had borrowed Fourier analysis from engineers; Gabor borrowed quantum
mechanics back. He pointed out that a signal cannot be both sharply located in time and sharply
located in frequency, proved the inequality~\eqref{eq:ch02:uncertainty} for signals, and showed
that the Gaussian-windowed sinusoid attains it. He proposed dividing the time--frequency plane into
cells of minimum area, which he called \emph{logons}, each able to carry one independent datum: a
signal of duration $T$ and bandwidth $B$ has roughly $2BT$ degrees of freedom. The same paper
introduced the \emph{analytic signal} (Section~\ref{sec:ch02:hilbert}) as the right way to define
the envelope and instantaneous frequency of a real signal. Gabor's expansion is the ancestor of the
short-time Fourier transform (Section~\ref{sec:ch02:stft}), of wavelets and of the time--frequency
grids on which OFDM and the delay--Doppler waveforms of Chapter~\ref{ch:ofdm} place their symbols.
Gabor went on to invent holography, for which he received the 1971 Nobel Prize in Physics.
\end{history}

%======================================================================
\section{Energy, power and correlation}
\label{sec:ch02:correlation}
%======================================================================

How does a GPS receiver find a signal that arrives far below the noise? How does a radar know
exactly when its echo came back? How does your Wi-Fi card know that a packet has just started? The
answer in every case is the same humble operation: slide a known pattern along what was received
and see where it fits best. That operation is correlation.

\mdcfig[0.85]{ch02_autocorr}{Left: the 13-chip Barker code. Centre: its aperiodic autocorrelation has a
peak of 13 and sidelobes no larger than 1 (22.3\,dB down), far better than a typical random
$\pm1$ sequence of the same length. Right: the Barker code's energy spectral density is nearly
13 times that of a single chip: the code is almost spectrally white within the chip bandwidth.}

\subsection{Energy spectral density and autocorrelation}

Parseval's theorem~\eqref{eq:ch02:parseval} says that $|X(f)|^2$ tells how the energy of a signal
is distributed in frequency: the energy between $f_1$ and $f_2$ is $\int_{f_1}^{f_2}|X(f)|^2df$.
The function $\Psi_x(f)=|X(f)|^2$ is the \term{energy spectral density} (ESD). It discards the
phase of $X(f)$, so it cannot distinguish signals that differ only in phase, but it is exactly what a
bandwidth or an interference calculation needs.

The time-domain partner of the ESD is the \term{autocorrelation function}
\begin{equation}
R_x(\tau)=\int_{-\infty}^{\infty}x(t+\tau)\,x^*(t)\,dt=\langle x(\cdot+\tau),x\rangle,
\label{eq:ch02:autocorr}
\end{equation}
the inner product of the signal with a shifted copy of itself. It measures how similar a signal
is to its own delayed version. Its transform is the ESD,
\begin{equation}
R_x(\tau)\ \leftrightarrow\ |X(f)|^2 ,
\label{eq:ch02:wkdet}
\end{equation}
which follows from the convolution and conjugation properties because
$R_x(\tau)=x(\tau)*x^*(-\tau)$. Several properties are immediate. $R_x(0)=E_x$; $|R_x(\tau)|\le R_x(0)$
by Cauchy--Schwarz, so the autocorrelation peaks at zero lag; and $R_x(-\tau)=R_x^*(\tau)$. A
signal with a wide spectrum has a narrow autocorrelation, and vice versa: the autocorrelation of a
rectangular pulse of width $T$ is a triangle of base $2T$, whose transform is the $T^2\sinc^2(fT)$
ESD. The \term{cross-correlation} $R_{xy}(\tau)=\int x(t+\tau)y^*(t)\,dt$ measures the similarity of
two different signals at every relative delay, and its transform is the cross-spectrum
$X(f)Y^*(f)$.

Correlation is how receivers find things (Figure~\ref{fig:ch02_correlation_find}). Synchronisation
(Chapter~\ref{ch:sync}) correlates the received signal with a known preamble and looks for the peak;
GPS receivers correlate with each satellite's spreading code to measure delay to a fraction of a chip
(Chapter~\ref{ch:spreadspectrum}); radar correlates the echo with the transmitted pulse. In every case
the ideal sequence has an autocorrelation like an impulse---one tall peak and low sidelobes---so that
the peak is unambiguous. By~\eqref{eq:ch02:wkdet} this is the same as asking for a flat ESD.

\begin{inpractice}[Barker codes, pulse compression and 802.11b]
A \term{Barker code} is a binary sequence whose aperiodic autocorrelation sidelobes are all 0 or
$\pm1$. Only a few exist; the longest known has length 13, giving a peak-to-sidelobe ratio of
$20\log_{10}13=22.3$\,dB (Figure~\ref{fig:ch02_autocorr}). Radars use them for \emph{pulse
compression}: transmit a long pulse whose phase flips according to the code, so that the energy
(and detection range) is that of the long pulse, then correlate the echo with the code, so that the
range resolution is that of a single chip. A 13\,$\mu$s pulse with 1\,$\mu$s chips gives the energy
of 13\,$\mu$s and the 150\,m resolution of 1\,$\mu$s---the escape from the trade-off of the radar
worked example above. Modern radars use linear-frequency-modulated chirps, whose compression ratios
reach the thousands. The original IEEE 802.11 and 802.11b Wi-Fi standards used the 11-chip Barker
sequence at 11\,Mchip/s for their 1 and 2\,Mb/s rates: the same correlation trick made the
signal robust to narrowband interference and multipath (Chapter~\ref{ch:spreadspectrum}).
\end{inpractice}

\mdcpairany{ch02_correlation_find}{Finding a needle in a haystack. A 13-chip Barker preamble (shaded)
is buried in noise stronger than itself; by eye it is invisible. Correlating with the known pattern
adds up all 104 samples coherently and produces an unmistakable peak at the right
place.}{ch02_pam_psd}{Random data inherits the spectrum of its pulse. Simulated power spectra of
random $\pm1$ symbols (wide pale curves) lie exactly on the energy spectrum $|P(f)|^2$ of a single
pulse (dashed), as~\eqref{eq:ch02:pampsd} predicts.}

\subsection{Power signals and the power spectral density}

For power signals the energy is infinite, so we average over time. The time-averaged
autocorrelation is
\begin{equation}
R_x(\tau)=\lim_{T\to\infty}\frac{1}{T}\int_{-T/2}^{T/2}x(t+\tau)\,x^*(t)\,dt ,
\end{equation}
with $R_x(0)=P_x$, and its Fourier transform is the \term{power spectral density} (PSD) $S_x(f)$,
the distribution of power over frequency, in watts per hertz (or dBm/Hz). The power in a band is the
integral of the PSD over the band, and the total power is $\int S_x(f)\,df=R_x(0)$. For a periodic
signal, the autocorrelation is periodic too and the PSD is a set of lines,
\begin{equation}
S_x(f)=\sum_k|c_k|^2\,\delta(f-kf_0),
\end{equation}
whose total weight is the Parseval sum~\eqref{eq:ch02:parsevalfs}.

For random signals---every data stream and every noise source---the time average is replaced by
an ensemble average, $R_x(\tau)=\E[x(t+\tau)x^*(t)]$, and the statement that the PSD is the
transform of the autocorrelation,
\begin{equation}
S_x(f)=\int_{-\infty}^{\infty}R_x(\tau)\,e^{-j2\pi f\tau}\,d\tau,
\label{eq:ch02:wk}
\end{equation}
is the \term{Wiener--Khinchin theorem}, named after Norbert Wiener (1930) and Aleksandr Khinchin
(1934). Chapter~\ref{ch:noise} develops it carefully, and Chapter~\ref{ch:baseband} uses it to
compute the spectrum of every line code. The PSD is what a spectrum analyzer displays, what a
regulator's emission mask constrains, and what we compute for every modulation format in this
book. One result is worth previewing because it ties this chapter together: a stream of
independent, zero-mean symbols $a_k$ of variance $\sigma_a^2$, each carried by a pulse $p(t)$ at
rate $1/T$, has PSD
\begin{equation}
S(f)=\frac{\sigma_a^2}{T}\,|P(f)|^2 .
\label{eq:ch02:pampsd}
\end{equation}
The spectrum of a random data signal has the \emph{shape} of the energy spectrum of a single pulse
(Figure~\ref{fig:ch02_pam_psd}). Every question about the bandwidth of a digital signal is therefore
a question about the Fourier transform of its pulse, which is why this chapter has spent so long on
pulses.

%======================================================================
\section{Linear time-invariant systems}
\label{sec:lti}
%======================================================================

Clap your hands once in a cathedral and listen. You hear the clap, then a rush of echoes, then a
long, fading roar as the sound bounces around the stone for several seconds. That roar is a complete
description of the building's acoustics: play any music in the cathedral and what reaches your ears
is the music ``smeared'' by exactly that echo pattern. This section turns that observation into the
most useful tool in engineering.

\mdcfig[0.85]{ch02_lti_rules}{The two rules that make a system LTI. Top: the response to a sum of
inputs is the sum of the individual responses (dotted). Bottom: delaying the input simply delays the
output. Everything in this section follows from these two pictures.}

A \term{system} turns an input signal into an output signal. Cables, filters, amplifiers,
antennas, multipath radio channels and whole receivers are systems. Two properties make a system
tractable. It is \term{linear} if the response to $a\,x_1+b\,x_2$ is $a\,y_1+b\,y_2$ for all inputs
and constants: superposition holds---two violins together sound like each violin alone, added. It is
\term{time-invariant} if delaying the input by $\tau$ simply delays the output by $\tau$: the system
behaves the same today as tomorrow---the cathedral's echo does not depend on what time the concert
starts. Wires, cables, passive filters, amplifiers operated well below compression, and radio
channels over intervals short compared with the time it takes the environment to change are all
well modelled as \emph{linear time-invariant} (LTI) systems. A squaring circuit is not linear; a
mixer, which multiplies by a local oscillator, is linear but not time-invariant; a fading channel is
linear but only approximately time-invariant over a short enough window (Chapter~\ref{ch:channels}).

\subsection{The impulse response and convolution}

An LTI system is completely described by one signal: its response $h(t)$ to an impulse
$\delta(t)$, the \term{impulse response}---the cathedral's answer to a single clap. The argument
is short and worth knowing. By the sifting property, any input is a superposition of shifted
impulses, $x(t)=\int x(\tau)\,\delta(t-\tau)\,d\tau$: a continuum of claps at times $\tau$ with
loudness $x(\tau)$. Time invariance says that the response to $\delta(t-\tau)$ is $h(t-\tau)$;
linearity says that the response to the weighted superposition is the same superposition of
responses. So
\begin{equation}
y(t)=\int_{-\infty}^{\infty}x(\tau)\,h(t-\tau)\,d\tau
=\int_{-\infty}^{\infty}h(\tau)\,x(t-\tau)\,d\tau
\quad\Longleftrightarrow\quad Y(f)=H(f)\,X(f).
\label{eq:ch02:conv}
\end{equation}
The integral is the \term{convolution} $y=x*h$, and its frequency-domain form follows from the
convolution property. Each output value is a weighted average of past and future inputs, the
weights being the impulse response read backwards.

\mdcphoto[0.6]{ch02_concert_hall}{The Golden Hall of the Vienna Musikverein, opened in 1870 and still
ranked among the world's finest-sounding rooms. Its tall, narrow ``shoebox'' shape and richly
decorated surfaces give it an impulse response---early reflections and a reverberant tail lasting
about two seconds---that flatters orchestral music.}

\mdcfig{ch02_room_echo}{Convolution as smearing. Four clean notes (top) pass through a room whose
impulse response (centre) has a direct sound, two early echoes and a decaying reverberant tail. The
output (bottom) is a scaled, delayed copy of $h(t)$ for each note, all added together.}

\begin{analogy}[Smearing each note with the room's echo]
Play four notes in a concert hall. Each note arrives at your seat followed by its own copy of the
hall's echo pattern, scaled by how loudly it was played and starting when it was played; what you
hear is all those smeared copies added together (Figure~\ref{fig:ch02_room_echo}). That sentence
\emph{is} the convolution integral: ``smear each input value with $h$, then add up''. The limit of
the analogy: a real hall is very nearly but not exactly LTI---the air warms up during a concert, the
audience shifts---and a radio channel can change in milliseconds as a car drives past.
\end{analogy}

Figure~\ref{fig:ch02_convolution} shows the recipe step by step: \emph{flip} $h(\tau)$ to get
$h(-\tau)$; \emph{slide} it to $h(t-\tau)$; \emph{multiply} by $x(\tau)$; \emph{integrate} the
product to get $y(t)$. Repeat for every $t$. With a rectangular input and the exponential impulse
response of an $RC$ low-pass filter, the overlap grows while the flipped exponential enters the pulse
(the output charges up), and shrinks once it has passed (the output discharges). The result is the
familiar smoothed pulse---and also, in miniature, the intersymbol interference of
Chapter~\ref{ch:baseband}: the tail of each smoothed pulse leaks into the time slots of the pulses
after it.

\mdcfig{ch02_convolution}{Convolution as flip, slide, multiply and integrate. A rectangular pulse
$x(\tau)$ is filtered by an $RC$ impulse response $h(t)=2e^{-t/0.5}u(t)$. Three panels show the
flipped and shifted response $h(t-\tau)$ at three times; the shaded area is the output at that
time. The fourth panel assembles the whole output.}

\begin{tryit}[Convolution in slow motion]
Lab~35's \emph{Convolution in slow motion} animates exactly this figure. Choose a
\emph{Rectangular pulse} and the \emph{RC low-pass}, then drag the \emph{Time t} slider and watch the
shaded overlap area build the output point by point. Switch the system to \emph{Echo (two paths)}:
the output becomes two copies of the input---the simplest multipath channel.
\end{tryit}

Convolution is commutative ($x*h=h*x$), associative ($(x*h_1)*h_2=x*(h_1*h_2)$) and
distributive over addition. Associativity means that a cascade of LTI systems is one LTI system
whose impulse response is the convolution of the individual responses and whose frequency
response is the product: transmit filter, channel and receive filter in
Chapter~\ref{ch:baseband} combine into a single end-to-end pulse. Commutativity means that the
order of LTI blocks does not matter---a fact used constantly in multirate DSP
(Chapter~\ref{ch:dsp}), and one that fails as soon as anything nonlinear or time-varying is in the
chain.

\subsection{Causality and stability}

A system is \term{causal} if its output never depends on future inputs, which for an LTI system
means $h(t)=0$ for $t<0$: the cathedral cannot echo a clap before you clap. Every physical system
operating in real time is causal; a digital filter processing a recorded file need not be, and many
optimum filters (the matched filter, the Wiener filter, the ideal low-pass filter) are not, until
they are delayed. A system is \term{BIBO stable} (bounded input, bounded output) if every bounded
input produces a bounded output, which for an LTI system holds exactly when $\int|h(t)|\,dt<\infty$.
An integrator, $h(t)=u(t)$, is unstable: a constant input produces a ramp. An $RC$ filter is stable;
an oscillator is, by design, on the edge of instability---the microphone-too-close-to-the-speaker
howl is a system that has crossed that edge.

\subsection{Frequency response: sinusoids in, sinusoids out}

Here is the reason sinusoids rule this chapter. Complex exponentials are the \emph{eigenfunctions}
of LTI systems: feed in $e^{j2\pi ft}$ and, from~\eqref{eq:ch02:conv}, out comes
$\int h(\tau)e^{j2\pi f(t-\tau)}d\tau=H(f)\,e^{j2\pi ft}$, the same exponential multiplied by a
complex number. That number,
\begin{equation}
H(f)=\int_{-\infty}^{\infty}h(t)\,e^{-j2\pi ft}\,dt=|H(f)|\,e^{j\theta(f)},
\end{equation}
is the \term{frequency response}, the Fourier transform of the impulse response. A real sinusoid
at frequency $f$ passes through with its amplitude multiplied by $|H(f)|$ and its phase advanced
by $\theta(f)=\angle H(f)$. An LTI system cannot create new frequencies, only scale and shift the
ones that are there---a tuning fork pressed against a table makes the table louder, never a different
note. The two curves that summarise an LTI channel are therefore its \emph{magnitude response}
$|H(f)|$, usually plotted in decibels, and its \emph{phase response} $\theta(f)$---and a channel is
``good'' when both are well behaved across the signal's band.

\mdcfig{ch02_rc_bode}{The $RC$ worked example in pictures. Left: the magnitude (navy, left axis) is
flat, passes $-3$\,dB at $f_3=159$\,kHz and then falls 20\,dB per decade, while the phase (red,
right axis) slides from $0^\circ$ through $-45^\circ$ at $f_3$ towards $-90^\circ$. Right: the step
response rises from 10\% to 90\% in 2.2\,$\mu$s.}

\begin{worked}[The $RC$ low-pass filter, from every angle]
A series resistor $R$ and shunt capacitor $C$ form a low-pass filter with
$h(t)=\frac{1}{RC}e^{-t/RC}u(t)$ and
\[
H(f)=\frac{1}{1+j2\pi fRC}=\frac{1}{1+jf/f_{3}},\qquad f_3=\frac{1}{2\pi RC}.
\]
With $R=1$\,k$\Omega$ and $C=1$\,nF: $RC=1\,\mu$s and $f_3=159$\,kHz (Figure~\ref{fig:ch02_rc_bode}).
At $f_3$ the magnitude is $1/\sqrt2$ ($-3$\,dB) and the phase is $-45^\circ$; far above $f_3$ the
response falls at 20\,dB per decade and the phase approaches $-90^\circ$. The step response rises
with time constant $1\,\mu$s, rising from 10\% to 90\% in $2.2RC=2.2\,\mu$s, the familiar rule
$t_r\approx0.35/f_3$. The \emph{group delay} (defined below) at low frequencies is
$RC=1\,\mu$s. The \emph{noise-equivalent bandwidth} (Section~\ref{sec:bandwidth}) is
$\int_0^\infty|H(f)|^2df=\frac{\pi}{2}f_3=250$\,kHz. Combining it with the thermal noise of the
resistor, $4kTR$ volts squared per hertz (Chapter~\ref{ch:noise}), gives an output noise of
$4kTR\times\frac{1}{4RC}=kT/C$, independent of $R$: $4.0\times10^{-21}/10^{-9}=4\times10^{-12}\,\text{V}^2$, or
$2.0\,\mu$V RMS. This ``$kT/C$ noise'' sets the noise floor of every switched-capacitor
sample-and-hold circuit, and therefore of most ADCs (Chapter~\ref{ch:pcm}).
\end{worked}

\subsection{Distortionless transmission}
\label{sec:ch02:distortionless}

What would a perfect channel do? It could make the signal weaker, and it could make it late---no
one minds a quieter, slightly delayed copy---but it must not change its shape. A channel that only
scales and delays, $y(t)=a\,x(t-\tau)$, has $H(f)=a\,e^{-j2\pi f\tau}$: constant magnitude and
\emph{linear} phase. The output is a faithful copy of the input. Any departure from either condition
across the band occupied by the signal causes distortion. \term{Amplitude distortion} changes the
relative sizes of frequency components: a telephone line that attenuates high frequencies more than
low ones rounds off pulse edges. \term{Phase distortion} misaligns components in time: even with a
perfectly flat magnitude, a nonlinear phase smears a pulse, because its frequency components arrive
at different times.

\mdcfig[0.8]{ch02_chord_arrival}{Group delay as a corridor that sorts notes by pitch. Three tone
bursts at 100, 250 and 400\,Hz start together (top). The channel passes all three at full strength
but delays each by its own group delay, $\tau_g(f)=0.05+0.0009f$\,s, so they arrive 0.14, 0.28 and
0.41\,s later (bottom). For a data pulse, made of many frequencies at once, the same effect smears the
pulse across its neighbours.}

\subsection{Phase delay, group delay and dispersion}
\label{sec:ch02:groupdelay}

Strike a chord on a piano at the far end of a long, strange corridor in which low notes travel
faster than high ones. You would hear not a chord but an arpeggio: the bass first, then the middle,
then the treble (Figure~\ref{fig:ch02_chord_arrival}). The corridor has not changed any note's
loudness, yet it has wrecked the music. That is phase distortion, and the quantity that measures it
is the group delay.

Two kinds of delay are defined from the phase response. The \term{phase delay}
\begin{equation}
\tau_p(f)=-\frac{\theta(f)}{2\pi f}
\label{eq:ch02:phasedelay}
\end{equation}
is the delay experienced by a pure sinusoid at frequency $f$: the time shift of its zero
crossings. The \term{group delay}
\begin{equation}
\tau_g(f)=-\frac{1}{2\pi}\frac{d\theta(f)}{df}
\label{eq:ch02:groupdelay}
\end{equation}
is the delay experienced by the \emph{envelope} of a narrowband signal centred on $f$. In words: the
phase delay is set by the phase itself, the group delay by its \emph{slope}. For a pure delay the
phase is a straight line through the origin and the two are equal. Pass a narrowband signal
$a(t)\cos(2\pi f_ct)$ through the system and the output is
\begin{equation}
y(t)\approx|H(f_c)|\,a(t-\tau_g)\cos\big(2\pi f_c(t-\tau_p)\big).
\label{eq:ch02:gdpacket}
\end{equation}
The envelope, which carries the information, travels at the group delay; the carrier's phase
travels at the phase delay (Figure~\ref{fig:ch02_group_delay}, left). When they differ, the
information arrives at $\tau_g$, and $\tau_g$ is the delay that matters for data, for ranging and
for latency.

\begin{deeper}[Why the envelope travels at the group delay]
Over the narrow band near $f_c$ occupied by the signal, expand the phase to first order:
$\theta(f)\approx\theta(f_c)+\theta'(f_c)(f-f_c)=-2\pi f_c\tau_p-2\pi(f-f_c)\tau_g$. The input's
spectrum is $\tfrac12A(f-f_c)+\tfrac12A(f+f_c)$, where $A(f)$ is the transform of $a(t)$. Near $+f_c$
the system multiplies it by $|H(f_c)|e^{-j2\pi f_c\tau_p}e^{-j2\pi(f-f_c)\tau_g}$. The factor
$e^{-j2\pi(f-f_c)\tau_g}$ is, by the time-shift property, a delay of $\tau_g$ applied to the
envelope $a(t)$; the constant factor $e^{-j2\pi f_c\tau_p}$ is a phase rotation of the carrier, a
shift of its zero crossings by $\tau_p$. Doing the same near $-f_c$ and adding the two halves
gives~\eqref{eq:ch02:gdpacket}.
\end{deeper}

\emph{Dispersion} is group delay that varies with frequency. A pulse occupies a range of
frequencies, so if each part of its spectrum arrives at a different time the pulse spreads. The
simplest case is a group delay that changes linearly across the band, $\tau_g(f)=2\pi\beta f$,
which corresponds to the quadratic phase $\theta(f)=-2\pi^2\beta f^2$ of
Figure~\ref{fig:ch02_group_delay} (right): a short Gaussian pulse emerges broader and lower, with its
energy unchanged and its frequency swept from one end of the band to the other (a chirp). Optical
fibre, the ionosphere, long copper cables and every sharp filter are dispersive.

\mdcfig{ch02_group_delay}{Left: a Gaussian-enveloped carrier through a system whose group delay
(1.2\,s) differs from its phase delay (0.3\,s). The envelope (dashed) arrives at the group delay;
the carrier's zero crossings are displaced by the phase delay. Right: dispersion. A short pulse
through a quadratic-phase channel, whose group delay varies linearly with frequency, spreads in
time; the stronger the dispersion $\beta$, the wider and lower the output.}

\begin{analogy}[Runners on a track]
A marathon starts as a tight pack. Each runner keeps a steady pace, but the paces differ, so after
42\,km the pack has stretched into a line many minutes long. Nothing was lost---every runner
finishes---but the shape of the pack is gone. A pulse in a dispersive fibre is the pack; its
frequency components are the runners. The analogy's limit: runners can overtake each other, while in
a linear channel the components pass through one another without interacting at all.
\end{analogy}

\begin{worked}[Group and phase delay in the ionosphere: GPS]
The ionosphere is a plasma whose refractive index for radio waves is slightly less than 1 and
depends on frequency. To first order, a signal at frequency $f$ crossing a path with total electron
content TEC (electrons per square metre) suffers an extra \emph{group} delay, expressed as a
distance, of
\[
\Delta_g=\frac{40.3\,\mathrm{TEC}}{f^2}\ \text{metres},
\]
and its carrier \emph{phase} is \emph{advanced} by the same amount: $\Delta_p=-\Delta_g$. Group
and phase delay have opposite signs. For the GPS L1 signal at 1575.42\,MHz and a moderate
$\mathrm{TEC}=10^{17}$\,m$^{-2}$ (10 TEC units), $\Delta_g=40.3\times10^{17}/(1.57542\times10^9)^2=1.62$\,m.
The spreading code, which rides on the envelope, arrives 1.62\,m ``late''; the carrier phase arrives
1.62\,m ``early''. Receivers that smooth code measurements with carrier phase see the two diverge,
and dual-frequency receivers exploit the $1/f^2$ law: on L2 at 1227.6\,MHz the delay is
$(1575.42/1227.6)^2=1.65$ times larger, 2.67\,m, so comparing the two measurements removes the
ionospheric error almost entirely (Chapter~\ref{ch:spreadspectrum}).
\end{worked}

\mdcfig[0.55]{ch02_iono_delay}{The ionosphere's $1/f^2$ law at a total electron content of 10 TEC
units: the code (envelope) is delayed and the carrier phase advanced by the same distance. L2 and L5
suffer about 1.65 and 1.8 times the error at L1, which is how dual-frequency receivers measure and
remove it.}

\begin{worked}[Chromatic dispersion in a fibre link]
Standard single-mode fibre has a dispersion coefficient of about $D=17$\,ps per nanometre of
optical bandwidth per kilometre at 1550\,nm: the group delay changes by 17\,ps for every nanometre
of wavelength, every kilometre. A 10\,Gb/s NRZ signal on an externally modulated laser occupies
roughly 0.08\,nm of optical spectrum (10\,GHz at 1550\,nm). Over 80\,km the spread of arrival
times is about $17\times80\times0.08\approx110$\,ps, a little more than one 100\,ps bit period. This is
why 10\,Gb/s links without compensation are limited to roughly 80\,km, and why the 100--400\,Gb/s
coherent systems of Chapter~\ref{ch:wireline}, whose symbols are shorter still, undo dispersion
digitally with an all-pass equaliser whose phase is the exact negative of the fibre's quadratic
phase. A 40\,Gb/s signal has four times the bandwidth and quarter-length bits, so its dispersion
tolerance is sixteen times worse: dispersion limits scale with the \emph{square} of the bit rate.
\end{worked}

\begin{inpractice}[Group delay in audio, cable and SerDes]
Group delay specifications appear wherever pulse shape or timing matters. \emph{Audio}: a
loudspeaker crossover network delays the woofer's band relative to the tweeter's by up to a
millisecond or so near the crossover frequency; the ear tolerates surprisingly large group-delay
variations at low frequencies but is more sensitive at mid frequencies, which is why studio
monitors advertise phase-linear designs. \emph{Cable television and DOCSIS}: the diplex filters
that separate upstream from downstream in a cable plant have steep skirts and therefore large
group-delay variation near the band edges; cable modems carry adaptive equalisers largely to
remove it (Chapter~\ref{ch:wireline}). \emph{Satellite transponders}: the input and output
multiplexer filters on board a satellite are sharp channel filters, and their group-delay
variation across the transponder is specified in nanoseconds and often pre-compensated on the
ground (Chapter~\ref{ch:satellite}). \emph{SerDes}: a printed-circuit-board channel at tens of
gigahertz has loss that rises with frequency and a phase that departs from linear; its
pulse response has a long tail of post-cursors, which the feed-forward and decision-feedback
equalisers of Chapter~\ref{ch:baseband} cancel. In every case the measurement is made with a
vector network analyser (Figure~\ref{fig:ch02_nanovna}), which measures phase directly and
computes~\eqref{eq:ch02:groupdelay} by differentiating it over a small aperture.
\end{inpractice}

\mdcphoto[0.55]{ch02_nanovna}{A NanoVNA, a pocket-sized vector network analyzer, showing a device's
reflection and transmission responses. Instruments like this measure magnitude \emph{and} phase
against frequency, from which the group delay follows by differentiation; a few decades ago the same
measurement needed a rack of equipment.}

\begin{tryit}[Turn a chord into an arpeggio]
In Lab~35, open \emph{Phase delay, group delay, dispersion}. With \emph{Pure delays}, set
\emph{Group delay} to 1.2\,s and \emph{Phase delay} to 0.3\,s and watch the envelope and the zero
crossings part company. Then choose \emph{Dispersion (quadratic phase)} and raise $\beta$: the pulse
spreads and chirps. Finally compare the \emph{Chebyshev} and \emph{Bessel} low-pass filters at the
same order and cut-off, and see which one keeps the pulse in shape.
\end{tryit}

\subsection{Ideal filters are impossible}

Every beginner wants the perfect filter: pass everything below $W$ untouched, block everything above
it completely. The \emph{ideal low-pass filter}, $H(f)=\operatorname{rect}(f/2W)$, does exactly that.
Its impulse response is the sinc pulse $h(t)=2W\sinc(2Wt)$, which extends from $t=-\infty$ to
$+\infty$. It responds before the impulse arrives; it is non-causal---it would have to hear the clap
before you clap (Figure~\ref{fig:ch02_ideal_lpf}). Delaying it by $t_d$ and truncating the part
before $t=0$ gives a causal approximation whose quality improves as $t_d$ grows---which is exactly
how FIR filters are designed (Chapter~\ref{ch:dsp}), and why sharp filters always have long delays.

\mdcfig[0.55]{ch02_ideal_lpf}{The ideal low-pass filter's impulse response (grey) starts ringing before
$t=0$, in the shaded ``future''. A real filter must delay it (here by 3 time units) and chop off what
is still left before zero. The sharper the filter, the longer the wait.}

Truncation is not merely inconvenient; perfection is forbidden. The \term{Paley--Wiener criterion}
states that a magnitude response $|H(f)|$ of finite energy can belong to a causal filter if and
only if
\begin{equation}
\int_{-\infty}^{\infty}\frac{\big|\ln|H(f)|\big|}{1+f^2}\,df<\infty .
\label{eq:ch02:paleywiener}
\end{equation}
If $|H(f)|$ is zero over any band of nonzero width, the logarithm is $-\infty$ there and the integral
diverges. So a causal filter cannot have a perfectly zero stopband; it can only make the stopband
very small. Nor can its magnitude fall faster than exponentially: a response like
$e^{-f^2}$ (a Gaussian) violates the criterion too, which is why a true Gaussian filter is also only
ever approximated. Infinite attenuation at isolated frequencies (the notches of an elliptic filter)
is allowed, because a logarithmic singularity at a point is integrable.

\subsection{Minimum phase and all-pass systems}

The magnitude response does not determine the phase. Given $|H(f)|$, many causal, stable filters
share it, differing by \term{all-pass} factors with $|A(f)|=1$ that alter only the phase. Among them,
exactly one has the smallest possible phase lag at every frequency: the \term{minimum-phase} filter.
In discrete time it is the one with all its zeros inside the unit circle; in continuous time, all its
zeros in the left half-plane. Minimum-phase systems have three properties that matter to engineers.
Their phase is fixed by their magnitude through a Hilbert-transform relation (the Bode gain--phase
relation): with the conventions of Section~\ref{sec:ch02:hilbert}, $\theta(f)=-\mathcal H\{\ln|H|\}(f)$.
Their energy is concentrated as early as possible: of all impulse responses with the same
magnitude response, the minimum-phase one has the most energy in its first $n$ samples, for every
$n$ (Figure~\ref{fig:ch02_minphase}). And they have stable causal inverses, which is why equalisers
(Chapter~\ref{ch:equalization}) can undo a minimum-phase channel exactly but must approximate the
inverse of a non-minimum-phase one with delay. A wireless channel whose strongest path arrives after a
weaker one is non-minimum phase, which is one reason equalisers have ``pre-cursor'' taps.

\mdcfig[0.75]{ch02_minphase}{Two echo channels with identical magnitude responses. In the
minimum-phase one (navy) the strong tap comes first; in the maximum-phase one (red) the weak tap
leads. Magnitude alone cannot tell them apart; an equaliser certainly can.}

\subsection{Filters: the classic trade-offs}

Figure~\ref{fig:ch02_filters} compares four classic fifth-order analog low-pass filter
families, all designed for a 1\,kHz cutoff, and it is worth studying like a family portrait. The
\emph{Butterworth} filter is maximally flat in the passband: the polite one. The \emph{Chebyshev}
filter allows passband ripple in exchange for a steeper transition. The \emph{elliptic} (Cauer)
filter allows ripple in both bands and gives the steepest transition of all for a given order: the
aggressive one. The \emph{Bessel} filter gives up selectivity for nearly constant group delay: the
one that never distorts a pulse. The time-domain panel reveals the price of selectivity: the sharp
filters have large group-delay peaks near the band edge, and a rectangular pulse passing through
them rings. For data signals, where the pulse shape is the information, this ringing is intersymbol
interference; that is why data receivers often prefer linear-phase digital FIR filters
(Chapter~\ref{ch:dsp}), whose group delay is exactly constant.

\mdcfig{ch02_filters}{Four fifth-order low-pass filters with a 1\,kHz cutoff. Left: magnitude.
Centre: group delay. Right: response to a 3\,ms rectangular pulse. Steeper magnitude
responses bring larger group-delay variation and more ringing.}

\begin{worked}[Filter selection for an audio channel]
A voice channel must pass 300--3400\,Hz and reject a 4\,kHz pilot tone by 40\,dB. The
transition band is narrow: the stopband edge is only $4000/3400=1.18$ times the passband edge, less
than a quarter of an octave. With 0.5\,dB allowed passband variation, the standard design formulas
(\texttt{scipy.signal.buttord} and its siblings compute them) give the required orders as 35
(Butterworth), 11 (Chebyshev) and 6 (elliptic). The elliptic design is the cheapest
in components but has the worst group delay; for voice, where the ear is insensitive to moderate
phase distortion, it is the usual choice. For a data modem using the same channel, the group-delay
variation would have to be equalized (Chapter~\ref{ch:equalization}).
\end{worked}

%======================================================================
\section{The Hilbert transform and the analytic signal}
\label{sec:ch02:hilbert}
%======================================================================

Look at an AM radio signal on an oscilloscope and your eye instantly draws its ``envelope''---the
smooth outline that the carrier's peaks trace out. But what, precisely, is that outline? The carrier
touches it only once per cycle; between peaks the outline is something the eye imagines. Defining
the envelope exactly, for any signal, turns out to require a new tool: a filter that shifts every
frequency by a quarter of a cycle.

A real signal's spectrum has Hermitian symmetry, $X(-f)=X^*(f)$, so its negative-frequency half
is redundant. Discarding it turns out to be the key to defining the envelope, phase and
instantaneous frequency of any signal, and to the complex baseband representation of the next
section. The tool that does the discarding is the Hilbert transform.

\mdcfig[0.85]{ch02_hilbert}{Left: an AM signal $x(t)$, its Hilbert transform $\hat x(t)$, and the
magnitude of the analytic signal, which is the envelope. Right: the instantaneous frequency of a
linear chirp computed from the phase of its analytic signal; the small glitches at the ends are
edge effects of the finite-length transform.}

\subsection{The Hilbert transform}

The \term{Hilbert transform} $\hat x(t)=\mathcal H\{x\}(t)$ is the output of an LTI filter with
frequency response
\begin{equation}
H_{\mathcal H}(f)=-j\,\operatorname{sgn}(f)=\begin{cases}-j=e^{-j\pi/2},&f>0,\\ +j=e^{+j\pi/2},&f<0,\end{cases}
\qquad h_{\mathcal H}(t)=\frac{1}{\pi t}.
\label{eq:ch02:hilbert}
\end{equation}
In words: it leaves every magnitude unchanged and shifts the phase of every positive-frequency
component by $-90^\circ$ (and every negative-frequency component by $+90^\circ$, as it must for a real
output). Every cosine becomes a sine: $\mathcal H\{\cos\}=\sin$ and $\mathcal H\{\sin\}=-\cos$. Its
properties follow from the frequency response: applying it twice multiplies by
$(-j\operatorname{sgn}f)^2=-1$, so $\mathcal H\{\mathcal H\{x\}\}=-x$; it preserves energy,
$\|\hat x\|=\|x\|$; and a signal is orthogonal to its own Hilbert transform,
$\langle x,\hat x\rangle=0$, which is the general form of the orthogonality of cosine and sine. The
impulse response $1/(\pi t)$ is non-causal and decays only as $1/t$, so a practical Hilbert
transformer is an approximation: an FIR filter with antisymmetric taps (Chapter~\ref{ch:dsp}) that
works over a band excluding DC, or a pair of all-pass networks whose phase responses differ by
$90^\circ$ across the band, the classic analog approach in SSB transmitters.

\subsection{The analytic signal: envelope and instantaneous frequency}

Now combine the signal with its quarter-cycle-shifted twin. The \term{analytic signal} of a real
signal $x(t)$ is
\begin{equation}
x_+(t)=x(t)+j\hat{x}(t),\qquad X_+(f)=X(f)+j\big(-j\operatorname{sgn}f\big)X(f)=2X(f)u(f),
\label{eq:ch02:analytic}
\end{equation}
where $u(f)$ is the unit step. Its spectrum is zero at negative frequencies and twice the original
at positive frequencies (Figure~\ref{fig:ch02_hilbert_spectra}); the original is recovered as
$x=\re\{x_+\}$. For a cosine, $x_+=\cos+j\sin=e^{j2\pi f_0t}$: the analytic signal of a sinusoid is
the single rotating phasor, with no counter-rotating partner---the bicycle wheel seen from the front.
That makes the polar form of $x_+$ a natural definition of a signal's \term{envelope} and
\term{instantaneous phase},
\begin{equation}
x_+(t)=a(t)\,e^{j\psi(t)},\qquad a(t)=|x_+(t)|,\qquad f_i(t)=\frac{1}{2\pi}\frac{d\psi(t)}{dt},
\label{eq:ch02:envphase}
\end{equation}
where $f_i$ is the \term{instantaneous frequency}: how fast the wheel is turning right now.
Figure~\ref{fig:ch02_hilbert} applies these definitions to an amplitude-modulated tone, where
$|x_+|$ traces the modulation exactly, and to a linear chirp, where the derivative of the analytic
signal's phase recovers the frequency sweep.

\mdcfig{ch02_hilbert_spectra}{How the analytic signal throws away negative frequencies. A real
signal's spectrum (left) is symmetric. Multiplying the Hilbert transform by $j$ gives a copy whose
negative-frequency half is flipped upside down (centre). Add the two and the negative half cancels,
while the positive half doubles (right).}

The envelope defined this way agrees with intuition when the signal is narrowband. For
$x(t)=a(t)\cos(2\pi f_ct+\phi(t))$ with $a(t)$ and $\phi(t)$ varying slowly enough that the
spectrum of $a(t)e^{j\phi(t)}$ lies within $|f|<f_c$, \emph{Bedrosian's theorem} gives
$\hat x(t)=a(t)\sin(2\pi f_ct+\phi(t))$ exactly, so that $|x_+|=a$ and the instantaneous
frequency is $f_c+\phi'(t)/2\pi$. When the signal is not narrowband these quantities still exist, but
they no longer mean ``the amplitude and frequency of a carrier''. Envelope detectors
(Chapter~\ref{ch:analog}), FM discriminators and the instantaneous-frequency estimators used in
frequency-hopping and radar receivers all compute, in effect, the polar form of an analytic
signal. Single-sideband modulation (Chapter~\ref{ch:analog}) is the analytic signal put to work:
$m(t)\cos(2\pi f_ct)\mp\hat m(t)\sin(2\pi f_ct)$ is $\re\{(m\pm j\hat m)e^{j2\pi f_ct}\}$, a carrier
modulated by an analytic (or anti-analytic) signal, with one sideband only.

%======================================================================
\section{Bandpass signals and the complex envelope}
\label{sec:complexenvelope}
%======================================================================

A Wi-Fi signal at 5.5\,GHz wiggles five and a half billion times a second. Yet the information it
carries changes only about twenty million times a second. If we had to simulate, filter and decode
every one of those carrier wiggles, a laptop would melt. We do not, because of the idea in this
section, which is the most important single idea in practical radio engineering: the carrier carries
no information, so take it away and describe only what is left.

\begin{analogy}[Describing a car relative to the highway]
Describe a car's motion relative to the Earth and you must include the Earth's spin, its orbit round
the Sun and the Sun's trip round the galaxy---huge, fast and irrelevant. Describe it relative to the
highway and you get what matters: it drifted half a metre left and slowed down a little. The complex
envelope is the highway frame for radio signals. The carrier is the Earth's spin, the same for every
signal on that frequency; the envelope is the lane change, and that is where the message lives
(Figure~\ref{fig:ch02_highway}). The analogy's limit: the highway frame is a convenience for the car,
but for radios it is more---it lets a single piece of DSP serve every carrier frequency at once.
\end{analogy}

Radio signals occupy a narrow band around a carrier frequency $f_c$ that is much larger
than their bandwidth: a 20\,MHz Wi-Fi channel at 5.5\,GHz, a 200\,kHz FM station at
100\,MHz, a 36\,MHz satellite transponder at 12\,GHz. Such \term{bandpass signals} can be
described completely by a low-frequency complex signal, the \term{complex envelope} or
\emph{baseband equivalent}. The rest of the book is written in it.

\mdcfig[0.8]{ch02_highway}{Two frames of reference for one QPSK signal. Left: the analytic signal
in the Earth frame whirls around the origin at the carrier frequency, tracing a dense tangle. Right:
in the highway frame---multiply by $e^{-j2\pi f_ct}$---the same signal is a slow, smooth path
whose positions at the symbol instants (dots) sit on the four QPSK points. All the information is in
the right-hand picture.}

\subsection{Three routes to the complex envelope}

There are three ways to arrive at the complex envelope, and seeing all three is what makes it
feel inevitable rather than clever.

\paragraph{From the physics.} Any bandpass signal can be written as a carrier whose amplitude and
phase vary slowly: $x(t)=a(t)\cos\big(2\pi f_ct+\phi(t)\big)$. Using Euler's formula,
\begin{equation}
x(t)=\re\big\{a(t)e^{j\phi(t)}\,e^{j2\pi f_ct}\big\}.
\end{equation}
All the information is in the slowly varying complex number $\tilde x(t)=a(t)e^{j\phi(t)}$; the
factor $e^{j2\pi f_ct}$ is a fixed, known rotation that carries no information.

\mdcfig{ch02_analytic}{From a real bandpass spectrum (left) to the analytic signal (centre)
to the complex envelope (right). The complex envelope need not be symmetric about zero.}

\paragraph{From the I/Q decomposition.} Expanding the cosine of a sum gives
\begin{equation}
x(t)=x_I(t)\cos(2\pi f_c t)-x_Q(t)\sin(2\pi f_c t),
\label{eq:ch02:iq}
\end{equation}
with $x_I=a\cos\phi$ and $x_Q=a\sin\phi$. The \term{in-phase} component $x_I$ modulates a cosine
carrier and the \term{quadrature} component $x_Q$ modulates a sine carrier. Because the two
carriers are orthogonal, both components fit into the same band and the receiver can separate them.
The complex envelope is simply $\tilde x=x_I+jx_Q$, the two real baseband signals packed into one
complex one.

\paragraph{From the spectrum.} Starting from the analytic signal~\eqref{eq:ch02:analytic}, which has
only the positive-frequency half, shift it down by $f_c$:
\begin{equation}
\tilde{x}(t)=x_+(t)\,e^{-j2\pi f_c t}=x_I(t)+j\,x_Q(t),\qquad
x(t)=\re\{\tilde{x}(t)\,e^{j2\pi f_c t}\}.
\label{eq:ch02:envelope}
\end{equation}
Figure~\ref{fig:ch02_analytic} shows the three spectra. The spectrum of the real signal is related
to that of its complex envelope by
\begin{equation}
X(f)=\tfrac12\tilde X(f-f_c)+\tfrac12\tilde X^*(-f-f_c),
\label{eq:ch02:envspec}
\end{equation}
a half-size copy at $+f_c$ and its mirror image at $-f_c$. The complex envelope is centred on zero
frequency and has the same bandwidth as the original signal's positive half, so it can be sampled
at a rate set by the signal bandwidth rather than the carrier frequency.

The last point in the caption deserves emphasis. A real baseband signal has a symmetric
magnitude spectrum; a complex envelope need not, because the part of the bandpass spectrum above
$f_c$ maps to positive frequencies and the part below $f_c$ maps to negative frequencies, and the two
are unrelated. An asymmetric complex envelope means that $x_I$ and $x_Q$ are correlated in a
particular way, and it is only possible because there are two real signals. A single-sideband
signal, a signal tuned slightly off-centre, and a band containing several unrelated channels all
have asymmetric complex envelopes.

Energy and power carry over with a factor of one half: with the convention~\eqref{eq:ch02:envelope},
$E_x=\tfrac12E_{\tilde x}$ and $P_x=\tfrac12P_{\tilde x}$, because averaging $\cos^2$ over a carrier
cycle gives $\tfrac12$. A constant envelope $\tilde x=A$ is the carrier $A\cos(2\pi f_ct)$, with
power $A^2/2$. Some books define $x=\sqrt2\re\{\tilde xe^{j2\pi f_ct}\}$ instead, so that powers are
equal in both domains; always check which convention a formula assumes.

\mdcfig{ch02_iq_helix}{Left: the complex exponential $e^{j2\pi ft}$ drawn against time is a
helix; its shadows on the two walls are a cosine (the real part, I) and a sine (the imaginary part,
Q). Right: an RRC-filtered QPSK signal. The real passband waveform (grey) oscillates at the carrier;
the complex envelope consists of the two slowly varying baseband signals $x_I$ and $x_Q$; its
magnitude is the envelope of the passband signal. Inset: the trajectory of the complex envelope in
the I/Q plane, with the symbol instants marked.}

\subsection{The I/Q modulator and demodulator}

Equation~\eqref{eq:ch02:iq} is a block diagram (Figure~\ref{fig:ch02_iqmod}). The
\term{I/Q modulator} multiplies $x_I$ by a cosine carrier and $x_Q$ by a negative sine
carrier---the same oscillator through a $90^\circ$ phase shifter---and adds them. The
\term{I/Q demodulator} reverses the process. Multiplying the received signal by
$2\cos(2\pi f_ct)$ and using the product formulas,
\begin{equation}
2x(t)\cos(2\pi f_ct)=x_I(t)+x_I(t)\cos(4\pi f_ct)-x_Q(t)\sin(4\pi f_ct),
\end{equation}
and a low-pass filter removes the terms at $2f_c$, leaving $x_I(t)$. Multiplying by
$-2\sin(2\pi f_ct)$ leaves $x_Q(t)$ in the same way. In complex notation the whole demodulator is
one line: multiply by $2e^{-j2\pi f_ct}$, which slides the positive-frequency copy of the spectrum to
zero and the negative-frequency copy to $-2f_c$, and low-pass filter.

\begin{figure}[htbp]\centering
\begin{tikzpicture}[node distance=0.5cm and 0.6cm,
  blk/.style={draw=navy,fill=navy!6,thick,minimum height=0.8cm,minimum width=1.1cm,align=center,font=\sffamily\scriptsize},
  mix/.style={draw=navy,circle,thick,minimum size=0.55cm,inner sep=0pt,font=\small},
  arr/.style={-{Stealth[length=2mm]},thick,navy}]
% modulator
\node[font=\sffamily\scriptsize] (xi) at (0,0.9) {$x_I(t)$};
\node[font=\sffamily\scriptsize] (xq) at (0,-0.9) {$x_Q(t)$};
\node[mix] (m1) at (1.5,0.9) {$\times$};
\node[mix] (m2) at (1.5,-0.9) {$\times$};
\node[blk] (lo) at (1.5,0) {LO $f_c$\\[-1pt]$0^\circ/{-90^\circ}$};
\node[mix,draw=accent] (sum) at (3.0,0) {$+$};
\draw[arr] (xi) -- (m1); \draw[arr] (xq) -- (m2);
\draw[arr] (lo) -- node[right,font=\tiny]{$\cos$} (m1);
\draw[arr] (lo) -- node[right,font=\tiny]{$-\sin$} (m2);
\draw[arr] (m1) -| (sum); \draw[arr] (m2) -| (sum);
\node[blk,fill=accent!6,draw=accent] (ch) at (4.6,0) {RF\\channel};
\draw[arr] (sum) -- node[above,font=\scriptsize]{$x(t)$} (ch);
% demodulator
\node[mix] (d1) at (6.3,0.9) {$\times$};
\node[mix] (d2) at (6.3,-0.9) {$\times$};
\node[blk] (lo2) at (6.3,0) {LO $f_c$\\[-1pt]$0^\circ/{-90^\circ}$};
\draw[arr] (ch) -- (5.5,0) |- (d1); \draw[arr] (5.5,0) |- (d2);
\draw[arr] (lo2) -- node[right,font=\tiny]{$2\cos$} (d1);
\draw[arr] (lo2) -- node[right,font=\tiny]{$-2\sin$} (d2);
\node[blk] (l1) at (7.9,0.9) {LPF};
\node[blk] (l2) at (7.9,-0.9) {LPF};
\draw[arr] (d1) -- (l1); \draw[arr] (d2) -- (l2);
\node[blk] (a1) at (9.4,0.9) {ADC};
\node[blk] (a2) at (9.4,-0.9) {ADC};
\draw[arr] (l1) -- (a1); \draw[arr] (l2) -- (a2);
\node[font=\sffamily\scriptsize,align=left] (o) at (11.0,0) {$\tilde x[n]=$\\$x_I[n]+jx_Q[n]$};
\draw[arr] (a1) -| (o); \draw[arr] (a2) -| (o);
\node[font=\sffamily\scriptsize\bfseries,text=navy] at (1.5,1.75) {I/Q modulator};
\node[font=\sffamily\scriptsize\bfseries,text=navy] at (7.9,1.75) {I/Q demodulator (zero-IF receiver)};
\end{tikzpicture}
\caption{The I/Q modulator and demodulator. Two real baseband signals ride on orthogonal carriers
through one RF channel. In the receiver, mixing with cosine and negative sine and low-pass filtering
recover them; two ADCs then deliver the complex envelope as a stream of complex samples. This is the
architecture of the AD9364 transceiver in the USRP B200 and of nearly every modern radio chip.}
\label{fig:ch02_iqmod}
\end{figure}

\begin{keyidea}[Why every SDR is an IQ machine]
A quadrature downconverter delivers $x_I$ and $x_Q$, and therefore the complete complex envelope,
at a sample rate set by the channel bandwidth rather than the carrier. One I/Q pair of ADCs at
61.44\,MS/s can capture any 56\,MHz-wide slice of the spectrum between 70\,MHz and 6\,GHz---that is
exactly what the USRP B200's AD9364 transceiver does. The same samples would represent the same
signal whatever the carrier, so all the DSP, from synchronisation to decoding, is carrier-independent.
And because complex sampling keeps positive and negative frequencies apart, the receiver can
distinguish signals above and below its tuning frequency, something a single real ADC at zero IF
cannot do. Chapter~\ref{ch:sdr} examines what happens when the I and Q paths are not perfectly
matched: the image of every signal leaks into its mirror frequency.
\end{keyidea}

\mdcphoto[0.55]{ch02_rtlsdr}{An inexpensive USB radio receiver of the kind sold for digital TV and
widely repurposed as a software-defined radio. A tuner chip converts a slice of spectrum down, and a
second chip digitises it and hands the computer a stream of complex I/Q samples---the complex envelope
of this section, delivered over USB.}

\begin{tryit}[Why a receiver needs two ADCs]
In Lab~1, open \emph{Why two ADCs: real vs IQ sampling}. Leave \emph{Tone 1 offset} at $+100$\,kHz
and \emph{Tone 2 offset} at $-250$\,kHz and compare the IQ spectrum with the I-only spectrum: with one
ADC every tone grows a mirror twin. Now drag tone 2 to $-100$\,kHz so it lands exactly on tone 1's
mirror, and see why a real-only zero-IF receiver cannot separate them.
\end{tryit}

\subsection{Baseband-equivalent systems}

If a bandpass signal passes through a bandpass system, the whole calculation can be done at
baseband, with complex numbers and a sample rate set by the bandwidth:
\begin{equation}
\tilde{y}(t)=\tilde{h}(t)*\tilde{x}(t)\qquad\Longleftrightarrow\qquad \tilde Y(f)=\tilde H(f)\,\tilde X(f).
\label{eq:ch02:bbequiv}
\end{equation}
The baseband-equivalent filter $\tilde H(f)$ is simply the positive-frequency part of the bandpass
response slid down to zero. This is why every simulation in this book, and in the labs, runs at
baseband: a channel at 28\,GHz is simulated with exactly the same code as one at 900\,MHz.

\mdcfig[0.8]{ch02_bb_equiv}{A bandpass filter tuned slightly above the carrier (left) has a
baseband equivalent that is lopsided about zero (right). A lopsided spectrum belongs to a complex
impulse response, which couples the I and Q channels.}

\begin{deeper}[Deriving the baseband-equivalent filter]
Write the system's impulse response as $h(t)=\re\{2\tilde{h}(t)e^{j2\pi f_c t}\}$, which in frequency
means $H(f)=\tilde H(f-f_c)+\tilde H^*(-f-f_c)$: the baseband-equivalent filter is the
positive-frequency part of the bandpass response slid down to zero, $\tilde H(f)=H(f+f_c)$ for
$f>-f_c$. Multiplying~\eqref{eq:ch02:envspec} by $H(f)$, the copy near $+f_c$ becomes
$\tfrac12\tilde H(f-f_c)\tilde X(f-f_c)$ (the cross terms vanish because the copies do not overlap),
which is the spectrum of a signal whose complex envelope is $\tilde h*\tilde x$.
\end{deeper}

A bandpass filter whose response is symmetric about $f_c$ has a real baseband equivalent and filters
$x_I$ and $x_Q$ independently. An asymmetric one---a filter tuned slightly off-centre, or a multipath
channel---has a complex $\tilde h$, and then the I output depends on both inputs: the filter couples I
into Q. That cross-coupling is what rotates and smears a constellation, and undoing it is what
equalisers do.

One consequence deserves special emphasis. A pure delay $\tau$ in the passband becomes, at
baseband, a delay $\tau$ \emph{and} a phase rotation:
\begin{equation}
x(t-\tau)\quad\Longrightarrow\quad\tilde x(t-\tau)\,e^{-j2\pi f_c\tau}.
\label{eq:ch02:delayrot}
\end{equation}
At a carrier of 2.4\,GHz, a path-length change of 12.5\,cm---one wavelength---rotates the phase by a
full cycle, while it changes the envelope delay by a negligible 0.4\,ns. That single fact explains
why carrier phase must be tracked (Chapter~\ref{ch:sync}), why multipath produces deep fades when
paths of nearly equal delay add with opposite phase (Chapter~\ref{ch:channels}), and how an antenna
array steers a beam by applying a different phase to each element (Chapter~\ref{ch:mimo}).

\mdcfig[0.5]{ch02_two_path_fade}{The worked example as a plot: received power as the receiver slides
along the line joining two paths (the second 0.8 times as strong). Peaks and 19\,dB nulls alternate
every 3\,cm---half a wavelength at 5\,GHz.}

\begin{worked}[Delay, phase and Doppler at baseband]
A 5\,GHz Wi-Fi signal reaches a receiver over two paths whose lengths differ by 3\,m. The extra
delay is $3/(3\times10^8)=10$\,ns, small compared with the 50\,ns sample period of a 20\,MHz channel,
so at baseband the two copies of the envelope nearly coincide. But the phase difference is
$2\pi f_c\tau=2\pi\times5\times10^9\times10^{-8}=100\pi$ radians, i.e., 50 whole cycles, so the two
copies arrive in phase and add. Move the receiver by 3\,cm, half a wavelength, along the line joining
the paths, and the phase difference changes by $\pi$: the copies cancel
(Figure~\ref{fig:ch02_two_path_fade}). If the receiver moves at a walking speed of 1.5\,m/s toward one
path, that path's length shrinks by 1.5\,m every second, which by \eqref{eq:ch02:delayrot} is a phase
that changes at $f_c v/c=25$\,Hz: a Doppler shift. Delay, phase rotation and Doppler are all the same
phenomenon seen through the complex envelope.
\end{worked}

\subsection{Sampling the complex envelope}

The complex envelope of a signal of bandwidth $B$ (total, two-sided around $f_c$) occupies
$-B/2<f<B/2$. Chapter~\ref{ch:pcm} shows that a complex signal confined to a band of width $B$ can be
reconstructed from samples at any rate $f_s\ge B$ complex samples per second---$2B$ real numbers per
second, exactly the same as a real signal of bandwidth $B$ sampled at its Nyquist rate. Nothing is
gained or lost by going complex; the information rate is the same. What is gained is convenience: the
carrier is gone, negative and positive frequencies are distinct, and every filter, mixer and detector
becomes complex arithmetic. A 20\,MHz LTE carrier is processed at 30.72\,MS/s complex; a 100\,MHz
5G NR carrier at 122.88\,MS/s; a 1\,MHz LoRa channel at a few megasamples per second. None of these
numbers depends on the carrier frequency.

\begin{pitfall}[Factors of two and the sign of $Q$]
Three conventions collide in the literature. (1) Whether $x=\re\{\tilde xe^{j2\pi f_ct}\}$ or
$x=\sqrt2\re\{\cdot\}$, which changes every power relation by a factor of two. (2) Whether the
quadrature carrier is $-\sin$ (this book, and~\eqref{eq:ch02:iq}) or $+\sin$, which conjugates the
complex envelope and mirrors its spectrum. (3) Whether the receiver mixes with $e^{-j2\pi f_ct}$ or
$e^{+j2\pi f_ct}$. Each choice is self-consistent, but mixing them produces mirrored spectra and
constellations that rotate the wrong way. When an SDR capture shows an FM station at $-200$\,kHz
that should be at $+200$\,kHz, swapping I and Q (or conjugating the samples) fixes it.
\end{pitfall}

%======================================================================
\section{Bandwidth}
\label{sec:bandwidth}
%======================================================================

Ask three engineers how wide a signal is and you may get three different answers, all correct. A
filter designer quotes the 3\,dB width; a regulator wants the band holding 99\% of the power; a
noise calculation needs the noise-equivalent width. For a carefully shaped signal they agree to
within a few per cent. For a careless one they can differ by a factor of twenty
(Figure~\ref{fig:ch02_occupied_bw})---the difference between passing and failing a regulatory test.

\mdcfig[0.62]{ch02_bandwidth}{Bandwidth definitions applied to the spectrum of a rectangular pulse
of duration $T$. Because the sidelobes decay only as $1/f^2$, the 99\% occupied bandwidth is about
ten times the null-to-null bandwidth.}

``Bandwidth'' sounds like a single number, but a signal with finite duration has a spectrum
that never quite reaches zero (Section~\ref{sec:ch02:uncertainty}), so ``the'' bandwidth has to be
defined by a convention. There are several useful conventions, each right for a particular job, and
confusing them causes real errors in system design. Figure~\ref{fig:ch02_bandwidth} illustrates the
classic ones for the spectrum of a rectangular pulse of duration $T$ (equivalently, the spectrum of
BPSK with rectangular pulses at symbol rate $1/T$, centred on its carrier).

\begin{description}[style=nextline,leftmargin=1.2em,font=\sffamily\bfseries\small\color{navy}]
\item[3\,dB (half-power) bandwidth.] The width of the band where the PSD is within 3\,dB of its peak;
for the rectangular pulse, $0.89/T$. The standard specification for filters and amplifiers, and a
reasonable description of where the bulk of a smooth spectrum lies.
\item[Null-to-null bandwidth.] The width of the main lobe, $2/T$ here. Convenient for spectra with
clear nulls (rectangular-pulse PSK, OFDM subcarriers), meaningless for spectra without them.
\item[Noise-equivalent bandwidth.] The width $B_N$ of an ideal rectangular filter with the same peak
gain that passes the same white-noise power as the actual filter,
\begin{equation}
B_N=\frac{\int_{-\infty}^{\infty}|H(f)|^2df}{\max_f|H(f)|^2}
\label{eq:ch02:noisebw}
\end{equation}
(two-sided; halve it for the one-sided convention common with real low-pass filters). It is the
right number to use in noise calculations (Chapter~\ref{ch:noise}). For a single-pole $RC$ filter
it is $\pi/2$ times the 3\,dB bandwidth; for a fifth-order Butterworth filter only about 2\% more
than the 3\,dB bandwidth; for a $\sinc^2$ spectrum exactly $1/T$; for a root-raised-cosine
filter exactly the symbol rate, whatever the roll-off.
\mdcfig[0.5]{ch02_occupied_bw}{The fraction of power captured as the bandwidth grows. A
root-raised-cosine signal reaches 99\% at $1.08/T$; a rectangular-pulse signal, dragging its slowly
decaying sidelobes, needs $20.5/T$.}

\item[Fractional-power (occupied) bandwidth.] The band containing a stated fraction of the total
power, usually 99\%, with 0.5\% left out on each side. This is the regulator's definition: the
International Telecommunication Union's Radio Regulations (No.~1.153) define \term{occupied
bandwidth} this way, and transmitter type approvals measure it. For the rectangular pulse the 99\%
bandwidth is about $20.5/T$, ten times the null-to-null width, because the sidelobes decay so slowly;
the 90\% bandwidth is only $1.7/T$.
\item[$x$-dB bandwidth.] The width between the outermost frequencies where the PSD falls $x$\,dB
below its peak (ITU-R Recommendation SM.328). The FCC uses the 26\,dB emission bandwidth for
unlicensed devices in the 5\,GHz U-NII bands.
\item[Necessary bandwidth.] The ITU's term (Radio Regulations No.~1.152) for the minimum bandwidth
sufficient for the emission's information rate and quality. It appears in the
\emph{emission designator} of every licence: in ``16K0F3E'' (narrowband FM voice), the 16K0 means a
necessary bandwidth of 16.0\,kHz, F frequency modulation, 3 a single analog channel and E telephony.
\item[Nyquist bandwidth.] The minimum bandwidth, $R_s/2$ at baseband or $R_s$ at passband, in which
symbols at rate $R_s$ can be sent without intersymbol interference (Chapter~\ref{ch:baseband}).
\end{description}

\subsection{Masks and adjacent-channel leakage}

Regulators and standards bodies rarely constrain a single bandwidth number. They specify a
\term{spectral emission mask}: a template of maximum PSD (in dBc per reference bandwidth, or dBm per
measurement bandwidth) as a function of offset from the carrier, which the transmitted spectrum
must stay under (Figure~\ref{fig:ch02_bandwidth_defs}). Think of it as the outline of a garage the
spectrum must park inside. Masks constrain the shoulders and skirts of the spectrum, where amplifier
nonlinearity (Section~\ref{sec:ch02:nonlinear}) and phase noise put power. A related figure of merit
is the \term{adjacent-channel leakage ratio} (ACLR), the ratio of the power in the assigned channel to
the power that falls in the neighbouring channel's band. The 3GPP specifications require, for
example, 45\,dB ACLR from LTE and 5G NR base stations (TS~36.104, TS~38.104) and about 30\,dB from
handsets (TS~36.101, TS~38.101-1). In practice it is the mask and the ACLR, not the occupied
bandwidth, that limit how hard a power amplifier can be driven.

\mdcfig[0.85]{ch02_bandwidth_defs}{One spectrum, five bandwidths: the PSD of root-raised-cosine QPSK
($\beta=0.22$) after a slightly nonlinear power amplifier, with the 3\,dB, noise-equivalent, 99\%
occupied and $-26$\,dB bandwidths marked, together with an illustrative spectral emission mask.
The shoulders at $\pm0.7$ to $\pm1.5R_s$ are spectral regrowth caused by the amplifier's third-order
nonlinearity.}

\begin{worked}[Two signals, five bandwidths]
Compare two 1\,Msymbol/s signals: BPSK with rectangular pulses, and QPSK with root-raised-cosine
pulses of roll-off $\beta=0.22$ (the WCDMA value) after the mildly nonlinear amplifier of
Figure~\ref{fig:ch02_bandwidth_defs}. From the analytic $\sinc^2$ spectrum and from the simulated PSD:

\begin{center}\small
\begin{tabular}{@{}lcc@{}}
\toprule
Definition & Rectangular BPSK & RRC QPSK, $\beta=0.22$\\
\midrule
3\,dB & 0.89\,MHz & 1.00\,MHz\\
Noise-equivalent & 1.00\,MHz & 0.99\,MHz\\
Null-to-null & 2.00\,MHz & (no nulls)\\
99\% occupied & 20.5\,MHz & 1.08\,MHz\\
$-26$\,dB & 11.3\,MHz & 1.22\,MHz\\
\bottomrule
\end{tabular}
\end{center}
The noise-equivalent bandwidths are identical: both receivers, using matched filters, collect the
same noise, and the two links have the same $E_b/N_0$ performance. Yet the rectangular signal
needs nearly twenty times the spectrum to satisfy a 99\% occupied-bandwidth rule. The 3\,dB
bandwidths differ by only 12\%. Quoting ``1\,MHz'' for both would be true by one definition and
wildly false by another, which is why every bandwidth number in a specification should name its
definition. (The $-26$\,dB figure for the rectangular pulse is set by its sidelobes: the
$\sinc^2$ sidelobe peaks, which fall as $1/(\pi fT)^2$, stay above $-26$\,dB out to about
$\pm5.65/T$.)
\end{worked}

\begin{tryit}[Five bandwidths of one signal]
Lab~35's \emph{Five bandwidths of one signal} marks all the definitions on a live spectrum. Start
with \emph{RRC-filtered QPSK} and drag the \emph{RRC roll-off} from 0.05 to 1: the 3\,dB and
noise-equivalent widths barely move while the 99\% width grows. Then add \emph{PA compression} and
watch the shoulders rise towards the mask. Finally select \emph{Rectangular BPSK} and see the 99\%
marker run off the screen.
\end{tryit}

\begin{pitfall}[Comparing bandwidths across definitions]
A datasheet may quote a filter's 3\,dB bandwidth, a standard may specify an occupied
bandwidth, and a noise calculation needs a noise-equivalent bandwidth. For a well-designed
root-raised-cosine signal the three are within 10--20\% of each other, which lulls engineers into
treating them as interchangeable; for a signal with sharp transitions they can differ by an order
of magnitude. Also check whether a quoted bandwidth is one-sided (baseband, $0$ to $W$) or two-sided
(passband, $f_c-W$ to $f_c+W$): a 1\,MHz baseband signal becomes a 2\,MHz passband signal after a DSB
mixer, and a complex baseband signal of ``bandwidth 20\,MHz'' spans $-10$ to $+10$\,MHz.
\end{pitfall}

%======================================================================
\section{Discrete time: from DTFT to FFT}
\label{sec:ch02:dft}
%======================================================================

Almost every spectrum you will ever look at---on a bench analyzer, in an SDR waterfall, in a
plot from your own Python script---was computed by a fast Fourier transform. The FFT is honest, but
it answers a slightly different question from the one you think you asked, and the difference can
hide a signal 60\,dB below a strong one, or misreport a tone's level by 4\,dB. This section develops
the discrete-time Fourier tools with an eye on the practical question: what does an FFT plot
actually show, and how can it mislead?

\mdcfig[0.55]{ch02_dtft_periodic}{Sampling makes the spectrum periodic: the same pattern repeats every
$f_s$. All the information sits in one period (dark), conventionally $-f_s/2\le f<f_s/2$.}

\subsection{The discrete-time Fourier transform}

If $x[n]=x(nT_s)$ is sampled at rate $f_s=1/T_s$, its spectrum is the
\term{discrete-time Fourier transform} (DTFT)
\begin{equation}
X(e^{j2\pi f/f_s})=\sum_{n=-\infty}^{\infty}x[n]\,e^{-j2\pi fn/f_s},
\label{eq:ch02:dtft}
\end{equation}
which is periodic in $f$ with period $f_s$. The periodicity is the frequency-domain face of
sampling---the wagon wheel again: since $e^{-j2\pi(f+f_s)n/f_s}=e^{-j2\pi fn/f_s}$ for every integer
$n$, a sampled signal cannot distinguish $f$ from $f+kf_s$. Equivalently
(Section~\ref{sec:ch02:delta}), sampling multiplies by an impulse train and therefore replicates the
spectrum every $f_s$. Chapter~\ref{ch:pcm} develops the sampling theorem and aliasing in full; for
now it is enough to remember that all the information of a correctly sampled signal lies in one
period, $-f_s/2\le f<f_s/2$, and that DSP engineers often measure frequency in cycles per sample,
$\nu=f/f_s$, or in radians per sample, $\omega=2\pi f/f_s$. All the properties of
Table~\ref{tab:ch02:props} have discrete-time counterparts, with sums in place of integrals.

\subsection{The discrete Fourier transform}

A computer can only handle a finite block of $N$ samples and a finite set of frequencies. The
\term{discrete Fourier transform} (DFT) evaluates the spectrum at $N$ equally spaced frequencies
$f_k=kf_s/N$, the \term{bins}:
\begin{equation}
X[k]=\sum_{n=0}^{N-1}x[n]\,e^{-j2\pi kn/N},\qquad
x[n]=\frac1N\sum_{k=0}^{N-1}X[k]\,e^{j2\pi kn/N}.
\label{eq:ch02:dft}
\end{equation}
\mdcfig[0.75]{ch02_dft_wrap}{What the DFT assumes. It treats the measured block (shaded) as one period
of an endlessly repeating signal. If the block holds a whole number of cycles (top), the repetition is
seamless and the tone falls in a single bin. If not (bottom), every block boundary has a jump, and a
jump---as Section~\ref{sec:ch02:fs} taught us---spreads energy over every harmonic.}

Three views of~\eqref{eq:ch02:dft} are useful. It is the DTFT of the $N$-sample block, sampled at $N$
frequencies. It is a change of basis: the vectors $e^{j2\pi kn/N}$, $k=0,\dots,N-1$, are orthogonal
over $N$ samples, and $X[k]$ is the projection of $x$ onto the $k$-th one---in matrix form
$\vect X=\mathbf F\vect x$, with $\mathbf F/\sqrt N$ a unitary matrix, and Parseval's theorem reads
$\sum|x[n]|^2=\frac1N\sum|X[k]|^2$. And it is the Fourier series of the periodic signal obtained by
repeating the block forever (Figure~\ref{fig:ch02_dft_wrap}). That implicit periodicity has three
consequences that dominate practical spectrum analysis: \emph{leakage}, the \emph{picket-fence
effect} and limited \emph{resolution}. It also means that multiplying two DFTs corresponds to
\emph{circular} convolution in time, the fact that OFDM's cyclic prefix is designed around
(Chapter~\ref{ch:ofdm}).

The bin spacing is $\Delta f=f_s/N$, the reciprocal of the observation time $NT_s$. For real input,
$X[N-k]=X^*[k]$ and only bins $0$ to $N/2$ are independent; for complex input all $N$ bins carry
information, with bins $N/2$ to $N-1$ representing negative frequencies (the \texttt{fftshift}
operation reorders them for display).

\subsection{Windows and leakage}
\label{sec:windows}

If a sinusoid does not complete an integer number of cycles in the block, the periodic
extension has a jump at the block boundary, and energy ``leaks'' from the true frequency into
all other bins, following the sidelobes of the sinc-shaped transform of the rectangular
block. Viewed through the multiplication property, the finite block is the infinite signal
multiplied by a rectangular \term{window}, so the computed spectrum is the true spectrum
convolved with the window's transform. Multiplying the block by a smooth window that tapers to zero
at both ends removes the jump and lowers the sidelobes---the Gibbs cure again.

\mdcfig[0.5]{ch02_window_frame}{Four windows seen in time: how quickly each fades the edges of the
block. The rectangular ``window'' does not fade at all; the flat-top fades fastest and even dips below
zero.}

\begin{analogy}[Looking through a picture frame]
Every FFT looks at the world through a frame. A plain rectangular frame has hard edges, and hard
edges are exactly what Fourier analysis hates: they ``leak'' light from a bright lamp all over the
picture. Fade the frame's edges smoothly to black and the glare disappears, at the price of seeing
the scene slightly less sharply (Figure~\ref{fig:ch02_window_frame}). The limit of the analogy:
a faded frame also dims what it shows, and a careful engineer must correct the amplitude and noise
readings for that dimming (the coherent gain and ENBW of Table~\ref{tab:ch02:windows}).
\end{analogy}

Figure~\ref{fig:ch02_windows} shows the trade. The rectangular window has the narrowest main lobe
but sidelobes only 13\,dB down, so a weak tone 60\,dB below a strong one 30\,Hz away is completely
hidden. The Hann window lowers sidelobes to $-31$\,dB with a faster roll-off, the Blackman--Harris
window to $-92$\,dB, and the flat-top window sacrifices resolution for amplitude accuracy.
Choosing a window is choosing between resolution, dynamic range and amplitude accuracy
(Table~\ref{tab:ch02:windows}).

\mdcfig{ch02_windows}{Left: DFT of a strong tone and a tone 60\,dB weaker, 30\,Hz away, with
four windows. Only low-sidelobe windows reveal the weak tone. Right: the spectra of the
windows themselves, in units of DFT bins.}

\begin{table}[htbp]\centering\small
\caption{Properties of common DFT windows (computed for long periodic windows; bandwidths in bins).
ENBW is the equivalent noise bandwidth; scalloping loss is the worst-case amplitude error for a tone
midway between bins.}\label{tab:ch02:windows}
\begin{tabular}{@{}lccccl@{}}
\toprule
Window & Sidelobe & 3\,dB width & ENBW & Scalloping & Typical use\\
\midrule
Rectangular & $-13.3$\,dB & 0.89 & 1.00 & 3.92\,dB & transients, OFDM\\
Hann & $-31.5$\,dB & 1.44 & 1.50 & 1.42\,dB & general purpose\\
Hamming & $-42.7$\,dB & 1.30 & 1.36 & 1.75\,dB & low close-in sidelobes\\
Kaiser ($\beta=9$) & $-66$\,dB & 1.67 & 1.76 & 1.07\,dB & adjustable trade-off\\
Blackman--Harris & $-92$\,dB & 1.90 & 2.00 & 0.83\,dB & high dynamic range\\
Flat-top & $-93$\,dB & 3.72 & 3.77 & 0.01\,dB & amplitude accuracy\\
\bottomrule
\end{tabular}
\end{table}

\begin{tryit}[Hide a tone, then find it]
In Lab~35's \emph{Windows, leakage and scalloping}, choose the \emph{Rectangular} window, set the
\emph{Weak tone level} to $-60$\,dB and the \emph{Distance} to 8 bins: the weak tone is invisible
under the strong tone's leakage. Now switch to \emph{Hann}, then \emph{Blackman--Harris}, and watch it
emerge. Slide \emph{Strong tone off-bin} between 0 and 0.5 to see both the leakage and the
scalloping loss change.
\end{tryit}

\subsection{The picket fence and scalloping loss}

The DFT samples the windowed spectrum only at the bin frequencies, like looking at a landscape
through a picket fence. A tone exactly on a bin is seen at its peak; a tone halfway between two bins
is seen on the shoulders of the window's main lobe, at a lower level
(Figure~\ref{fig:ch02_scalloping}). This \term{scalloping loss} is 3.92\,dB for the rectangular
window, 1.42\,dB for Hann, and essentially zero for a flat-top window, whose main lobe is
deliberately flat across a whole bin. That is why spectrum analyzers offer a flat-top window for
amplitude measurements, and why a tone's power measured from a single FFT bin can be in error by
several decibels unless the window is chosen accordingly. Summing the power over the few bins of the
main lobe, or interpolating the peak (fitting a parabola to the three highest bins in decibels is a
common trick), recovers both the true amplitude and a frequency estimate far finer than the bin
spacing.

\mdcfig[0.85]{ch02_scalloping}{Left: with a rectangular window, a tone on a bin centre appears in one bin,
while a tone halfway between bins appears 3.9\,dB low and leaks into every bin. Right: the scalloping
loss against tone offset for four windows.}

\subsection{Resolution and zero padding}

Two tones closer together than about $f_s/N$\,Hz---the reciprocal of the observation time
$NT_s$---cannot be resolved, no matter how the data are processed; with a window, the limit is the
window's main-lobe width, one to four bins. Padding the block with zeros computes the spectrum at more
closely spaced frequencies, which makes plots smoother and helps locate an isolated peak precisely,
but it does not add resolution (Figure~\ref{fig:ch02_dft_resolution}). It interpolates the same
underlying DTFT of the same windowed block. Resolution requires observation time. This is the
discrete form of the uncertainty principle: a measurement of duration $T$ cannot distinguish
frequencies closer than about $1/T$---the drum and the violin once more.

\mdcfig[0.85]{ch02_dft_resolution}{Left: two tones 1.2\,Hz apart are merged with 64 samples at
100\,Hz and resolved with 256. Right: zero padding a 32-point DFT interpolates the same
underlying spectrum more finely.}

\begin{pitfall}[Zero padding is not free resolution]
A beautifully smooth, zero-padded FFT plot looks high-resolution, and that is exactly the danger.
Padding a 1\,ms capture to a million points gives you a curve sampled every hertz, but two tones
closer than about 1\,kHz will still merge into one hump. If you need to separate them, capture for
longer.
\end{pitfall}

\begin{tryit}[Pad, then measure longer]
In Lab~35's \emph{Zero padding is not resolution}, set the \emph{Tone spacing} to 10\,Hz and the
\emph{Record length N} to 32: the two tones merge. Raise \emph{Zero padding} to $\times64$---the
curve becomes smooth but stays merged. Now increase $N$ to 256 instead, and the tones separate.
\end{tryit}

\begin{worked}[Choosing FFT parameters for a measurement]
An engineer wants to see whether a 2.4\,GHz transmitter has spurs 50\,dB below the carrier and
within 20\,kHz of it, using an SDR sampling at $f_s=1$\,MS/s complex. She needs resolution well
under 20\,kHz and sidelobes below $-50$\,dB. Choose a Blackman--Harris window (sidelobes $-92$\,dB,
main lobe $\pm4$ bins). To separate a spur 20\,kHz away clearly, the main lobe's half-width, 4 bins,
should be at most about 5\,kHz, so the bin spacing must be $\le1.25$\,kHz and $N\ge
10^6/1250=800$; take $N=1024$, a bin spacing of 977\,Hz and an observation time of 1.02\,ms. The
equivalent noise bandwidth per bin is $2.0\times977\approx2$\,kHz, so a noise density of
$-150$\,dBm/Hz appears at $-150+10\log_{10}2000=-117$\,dBm per bin. Averaging the power spectra of
many such blocks (Welch's method) reduces the scatter of the noise floor without changing its mean,
making weak spurs stand out (Figure~\ref{fig:ch02_welch}).
\end{worked}

\mdcpairany{ch02_welch}{Welch averaging in action, with the worked example's settings (1\,MS/s,
1024-point Blackman--Harris FFT). A single FFT of noise (grey) scatters over 30\,dB; the average of
200 (navy) is a smooth floor, and a weak spur at 123\,kHz stands clear of it.}{ch02_fft_ops}{Direct
DFT against FFT. At $N=4096$ the FFT needs 680 times fewer multiplications; at the 65\,536-point
sizes used in radio astronomy and spectrum monitoring the factor exceeds 8000.}

\subsection{The fast Fourier transform}

A direct DFT takes $N^2$ complex multiplications. For a 4096-point transform that is 16.8 million
multiplications---every 12.8\,$\mu$s, in a Wi-Fi~7 receiver. The \term{fast Fourier transform}
(FFT) computes exactly the same numbers with about $\frac N2\log_2N$, and it is the reason your phone
can do OFDM at all. The idea is divide and conquer: a big DFT is two half-size DFTs glued together
with a little extra arithmetic, and each half can be split again, and again.

\begin{deeper}[The radix-2 split]
For $N$ a power of two, split the sum~\eqref{eq:ch02:dft} into even- and odd-indexed samples.
Writing $W_N=e^{-j2\pi/N}$,
\begin{equation}
X[k]=\underbrace{\sum_{m=0}^{N/2-1}x[2m]\,W_{N/2}^{km}}_{E[k]}\;+\;W_N^k\underbrace{\sum_{m=0}^{N/2-1}x[2m+1]\,W_{N/2}^{km}}_{O[k]},
\label{eq:ch02:fftsplit}
\end{equation}
where $E$ and $O$ are $N/2$-point DFTs of the even and odd samples. Because $E$ and $O$ are periodic
with period $N/2$ and $W_N^{k+N/2}=-W_N^k$, the two halves of the output come from the same pair of
numbers:
\begin{equation}
X[k]=E[k]+W_N^kO[k],\qquad X[k+N/2]=E[k]-W_N^kO[k],\qquad k=0,\dots,N/2-1 .
\end{equation}
This two-input, two-output operation is the \term{butterfly}.
\end{deeper}

Each half-size DFT is split the same way, recursively, until only 2-point DFTs remain. There are
$\log_2N$ stages of $N/2$ butterflies, each needing one complex multiplication, for
$\frac N2\log_2N$ multiplications and $N\log_2N$ additions in total (Figure~\ref{fig:ch02_butterfly}).
For $N=4096$ that is 24\,576 multiplications instead of 16.8 million, a factor of about 680
(Figure~\ref{fig:ch02_fft_ops}). Splitting the data in this way scrambles the input order into
\emph{bit-reversed} order (sample 1, binary 001, goes to position 4, binary 100), which hardware FFTs
handle with a reordering memory.

\begin{figure}[htbp]\centering
\begin{tikzpicture}[x=1cm,y=0.5cm]
\foreach \i/\lab in {0/0,1/4,2/2,3/6,4/1,5/5,6/3,7/7} {
  \node[left,font=\scriptsize] at (0.3,-\i) {$x[\lab]$};
  \node[right,font=\scriptsize] at (7.2,-\i) {$X[\i]$};
  \foreach \c in {0.3,2.6,4.9,7.2} \fill[navy] (\c,-\i) circle (1.3pt);
}
\foreach \a/\b in {0/1,2/3,4/5,6/7} {
  \draw[navy] (0.3,-\a)--(2.6,-\a); \draw[navy] (0.3,-\b)--(2.6,-\b);
  \draw[navy] (0.3,-\a)--(2.6,-\b); \draw[navy] (0.3,-\b)--(2.6,-\a);
}
\foreach \a/\b/\k in {0/2/0,1/3/2,4/6/0,5/7/2} {
  \draw[navy] (2.6,-\a)--(4.9,-\a); \draw[navy] (2.6,-\b)--(4.9,-\b);
  \draw[navy] (2.6,-\a)--(4.9,-\b); \draw[navy] (2.6,-\b)--(4.9,-\a);
  \node[accent,font=\tiny,fill=white,inner sep=0.5pt] at (3.0,-\b+0.35) {$W_8^{\k}$};
}
\foreach \a/\b/\k in {0/4/0,1/5/1,2/6/2,3/7/3} {
  \draw[navy] (4.9,-\a)--(7.2,-\a); \draw[navy] (4.9,-\b)--(7.2,-\b);
  \draw[navy] (4.9,-\a)--(7.2,-\b); \draw[navy] (4.9,-\b)--(7.2,-\a);
  \node[accent,font=\tiny,fill=white,inner sep=0.5pt] at (5.3,-\b+0.35) {$W_8^{\k}$};
}
\node[font=\sffamily\scriptsize,text=navy] at (1.45,1.0) {stage 1: 2-point DFTs};
\node[font=\sffamily\scriptsize,text=navy] at (3.75,1.0) {stage 2: 4-point};
\node[font=\sffamily\scriptsize,text=navy] at (6.05,1.0) {stage 3: 8-point};
\end{tikzpicture}
\caption{Signal-flow graph of an 8-point radix-2 decimation-in-time FFT. Inputs are in bit-reversed
order; each stage combines pairs of smaller DFTs with butterflies. In every butterfly the lower input
is multiplied by the twiddle factor shown ($W_8=e^{-j2\pi/8}$; stage~1 twiddles are all 1), then
added to the upper input to form the upper output and subtracted to form the lower output. Three
stages of four butterflies replace 64 complex multiplications by 12, most of them trivial.}
\label{fig:ch02_butterfly}
\end{figure}

Production FFT libraries (FFTW, Intel MKL, cuFFT, the FFT engines in every OFDM modem chip) go well
beyond radix 2. They use radix-4 and split-radix decompositions, which save about a quarter of the
multiplications; mixed radices for lengths like $1536=3\times2^9$ (used in LTE); prime-factor and
Bluestein algorithms for awkward lengths; and cache-aware orderings that matter more on a modern CPU
than the arithmetic count. The FFT made digital spectrum analysis routine, and forty years after its
publication it made OFDM practical: a Wi-Fi~7 receiver computes a 4096-point FFT every
12.8\,$\mu$s (Chapter~\ref{ch:ofdm}).

\begin{history}[Cooley, Tukey---and Gauss]
In 1965 James Cooley of IBM and John Tukey of Princeton published ``An algorithm for the machine
calculation of complex Fourier series'' in \emph{Mathematics of Computation}, a five-page paper that
became one of the most cited in applied mathematics. The story usually told is that Tukey sketched
the idea at a meeting of President Kennedy's Science Advisory Committee, where Richard Garwin of IBM
was interested in fast spectral analysis---among other uses, for detecting nuclear tests in seismic
records---and pressed Cooley to program it. Its effect was immediate: computations that took hours took
seconds, and within a few years the FFT had spawned digital spectrum analyzers, fast convolution and
the field of digital signal processing. Then the historians went looking. Variants had been published
by Danielson and Lanczos (1942), by Good (1958) and by Runge in the early 1900s. And in 1984 Michael
Heideman, Don Johnson and Sidney Burrus showed that Carl Friedrich Gauss had used essentially the same
divide-and-conquer recursion in about 1805, to interpolate the orbits of the asteroids Pallas and
Juno from observations---before Fourier presented his memoir. Gauss's notes were written in Latin and
published only posthumously in his collected works, where the idea sat unnoticed for 160 years.
Good algorithms, it seems, are discovered when the computers are ready for them.
\end{history}

\mdcphoto[0.32]{ch02_gauss}{Carl Friedrich Gauss in 1840. Around 1805, while fitting the orbits of
asteroids by hand, he used the divide-and-conquer trick now called the FFT---and left it in an
unpublished Latin manuscript for 160 years.}

\begin{inpractice}[Real-signal FFT tricks]
Many signals are real---audio, the output of a single ADC, the real and imaginary parts of separate
measurements---and a complex FFT wastes half its effort on them, because the negative-frequency half
of the output is the conjugate of the positive half. Two classical tricks recover the waste.
\emph{Two real FFTs for the price of one}: pack two real $N$-sample signals as $z=x+jy$, compute
$Z=\mathrm{FFT}(z)$, and separate the results using the symmetries
$X[k]=\tfrac12\big(Z[k]+Z^*[N-k]\big)$ and $Y[k]=\tfrac1{2j}\big(Z[k]-Z^*[N-k]\big)$. \emph{One
$2N$-point real FFT from one $N$-point complex FFT}: put the even samples in the real part and the odd
samples in the imaginary part, transform, separate as above to get the even and odd DFTs, and combine
them with one more stage of butterflies~\eqref{eq:ch02:fftsplit}. Library routines such as
\texttt{numpy.fft.rfft} and FFTW's \texttt{r2c} transforms do this internally and return only the
$N/2+1$ non-redundant bins. Two further everyday tricks: to compute a correlation or a long
convolution, use $\mathrm{IFFT}(\mathrm{FFT}(x)\cdot\mathrm{FFT}(y)^*)$ with enough zero padding to
avoid circular wrap-around (Chapter~\ref{ch:dsp}); and to shift a spectrum by half the sample rate,
multiply the time samples by $(-1)^n$, which costs nothing but sign changes.
\end{inpractice}

\subsection{Spectrum analyzers, RBW and VBW}

Two families of instrument display spectra, and both are explained by this chapter. A classic
\emph{swept} spectrum analyzer is a superheterodyne receiver (Chapter~\ref{ch:analog}) whose local
oscillator sweeps across the span, like a radio whose tuning knob is turned steadily from one end of
the dial to the other while a pen records the volume. The signal at each frequency is converted to a
fixed intermediate frequency, filtered by the \term{resolution bandwidth} (RBW) filter, detected and
displayed. The RBW filter plays the role of the window in an FFT: it sets the resolution (two tones
closer than about one RBW merge) and the displayed noise floor, which rises by $10\log_{10}$ of the
RBW---a tenfold narrower RBW lowers the noise floor by 10\,dB and reveals weaker signals
(Figure~\ref{fig:ch02_rbw}). The price is time. The RBW filter needs about $1/\text{RBW}$ to respond,
and the sweep must dwell that long on each RBW-wide slice, so the sweep time grows roughly as
$k\cdot\text{span}/\text{RBW}^2$, with $k$ of order 2--3 for typical filters. A 1\,GHz span at
10\,kHz RBW takes a few tens of seconds; at 1\,kHz it takes nearly an hour. After the detector, a
\term{video bandwidth} (VBW) filter smooths the displayed trace; setting VBW well below RBW averages
the noise into a smooth line, at the cost of further slowing the sweep.

\mdcpairany{ch02_rbw}{One input---two tones 30\,kHz apart plus noise---viewed with three resolution
bandwidths. At 100\,kHz the tones merge into one hump; at 10\,kHz they separate; at 1\,kHz the noise
floor drops another 10\,dB and the tones are sharp spikes.}{ch02_led_analyzer}{A home-made audio
spectrum display built from an LED matrix and a microcontroller running an FFT: the same window,
bin and averaging choices as a laboratory analyzer, on a hobbyist's budget.}

An \emph{FFT} or \emph{real-time} analyzer instead digitises a whole band at once, typically up to a
few hundred megahertz, and computes overlapping windowed FFTs. Its RBW is the window's equivalent
noise bandwidth times the bin spacing, $\text{ENBW}\cdot f_s/N$. Because every frequency in the
band is observed simultaneously, it captures brief and intermittent signals that a swept analyzer
would miss, and with enough overlap between FFT frames it misses nothing at all---the ``real-time''
in the name. Modern bench analyzers combine both: they sweep a digital IF across wide spans and use
FFTs at narrow RBWs, where sweeping would be slow. GNU Radio's \texttt{QT GUI Frequency Sink} and the
waterfall display of every SDR program are FFT analyzers; their FFT size and window set the RBW.

\mdcphoto[0.55]{ch02_specan2}{The front panel of a classic swept spectrum analyzer. The
frequency, span and marker keys, and the resolution- and video-bandwidth settings discussed below,
are the controls every RF engineer reaches for first.}

\begin{inpractice}[Measuring noise and modulated signals on a spectrum analyzer]
A spectrum analyzer is calibrated to read the power of a \emph{sinusoid} correctly. Reading noise or
a noise-like digitally modulated signal needs corrections. First, the power displayed is the power
in one RBW, so a density follows by subtracting $10\log_{10}(\text{RBW}_{\text{noise}})$, where the
noise bandwidth of a typical Gaussian-shaped RBW filter is about 1.06 times its 3\,dB width: a noise
floor read as $-100$\,dBm in a 100\,kHz RBW is about $-100-50.3=-150.3$\,dBm/Hz. Second, a
log-scale display averages the logarithm of the noise envelope, which reads about 2.5\,dB low for
Gaussian noise; analyzers' ``noise marker'' functions apply both corrections automatically. Third,
the power of a wideband signal is the integral of its PSD over the channel, which the ``channel
power'' and ``ACLR'' measurement modes compute by summing bins with RMS detection. A 20\,MHz LTE
carrier displayed in a 30\,kHz RBW looks 28\,dB weaker than its total power, $10\log_{10}(18\,\text{MHz}/30\,\text{kHz})$,
because each RBW slice holds only a small fraction of it. Forgetting this is the commonest
spectrum-analyzer mistake in the laboratory.
\end{inpractice}

%======================================================================
\section{Time--frequency analysis: the short-time Fourier transform}
\label{sec:ch02:stft}
%======================================================================

A musical score is a time--frequency diagram: time runs left to right, pitch runs up the stave, and
each note is a mark at a place and a height. The Fourier transform, by contrast, is like a list of
every note in the symphony with no indication of when it was played. For a steady tone that is fine;
for anything that changes---speech, a frequency-hopping radio, a packet that starts and stops---we
want the score.

\mdcfig{ch02_spectrogram}{Spectrograms. Left and centre: a linear chirp (200\,Hz to 2.8\,kHz in
1\,s) plus a 60\,ms tone burst at 1.5\,kHz, analysed with 8\,ms and 128\,ms Hann windows. The short
window resolves time and blurs frequency; the long window does the opposite. Right: a 4-FSK burst
with 20\,ms symbols, which the spectrogram decodes visually.}

The Fourier transform describes a signal's frequency content over all time; it cannot say
\emph{when} a frequency occurred. A chirp and a click can have identical magnitude spectra. Real
signals change: a frequency-hopping radio jumps channel every few hundred microseconds, a speech
formant glides, a packet starts and stops. The \term{short-time Fourier transform} (STFT) slides a
window $w(t)$ along the signal and transforms each windowed segment,
\begin{equation}
X(\tau,f)=\int_{-\infty}^{\infty}x(t)\,w(t-\tau)\,e^{-j2\pi ft}\,dt ,
\label{eq:ch02:stft}
\end{equation}
and the \term{spectrogram} $|X(\tau,f)|^2$, drawn as an image with time across and frequency up (or
as a scrolling ``waterfall''), shows how the spectrum evolves. In discrete time it is simply a
sequence of windowed FFTs over overlapping blocks.

The window length sets a trade-off that is the uncertainty principle again. A short window locates
events precisely in time but blurs frequency (its spectral resolution is about the reciprocal of its
length); a long window resolves frequency finely but smears events in time. Figure~\ref{fig:ch02_spectrogram}
shows a linear chirp sweeping 2.6\,kHz per second, with a short tone burst. With an 8\,ms window the
burst's start and stop are sharp but every line is about 200\,Hz wide; with a 128\,ms window the
steady tone is a thin line, but the chirp, which sweeps 330\,Hz during one window, is smeared into a
broad stripe and the burst's edges blur. For a chirp of rate $\alpha$\,Hz/s the best compromise is a
window of length about $1/\sqrt\alpha$, here about 20\,ms, which balances the frequency blur $1/T$
against the sweep $\alpha T$ during the window. The third panel shows a 4-FSK burst, whose
symbol-by-symbol frequency hops are read off directly. Engineers use spectrograms constantly: to find
intermittent interference, to verify frequency-hopping patterns (Bluetooth hops 1600 times per second,
Chapter~\ref{ch:wifi}), to inspect LoRa chirps, and to see packet timing at a glance
(Figure~\ref{fig:ch02_ism_waterfall}).

\mdcfig[0.7]{ch02_tf_tiles}{Gabor's tiles. Every STFT divides the time--frequency plane into cells of
the same minimum area. A short window makes them tall and thin (sharp in time); a long window makes
them short and wide (sharp in frequency). No window makes them smaller.}

\mdcfig[0.65]{ch02_ism_waterfall}{A synthetic waterfall of the kind an SDR shows in a busy unlicensed
band: short narrowband bursts hopping across the band, two wideband packets, and a train of chirps. A
plain spectrum would mix all of these into one smear; the spectrogram separates them at a glance.}

Gabor's time--frequency cells (the history box of Section~\ref{sec:ch02:uncertainty}) are the tiles of
the STFT: no window can make them smaller than the uncertainty bound, only change their aspect ratio
(Figure~\ref{fig:ch02_tf_tiles}). Wavelet transforms change the aspect ratio with frequency, using
short windows at high frequencies and long ones at low frequencies, which suits signals like music
and images. OFDM can be seen as an STFT used backwards: the transmitter places one data symbol in each
time--frequency tile, and the receiver's FFT reads them out (Chapter~\ref{ch:ofdm}).

\begin{tryit}[Find the best window for a chirp]
In Lab~35's \emph{The spectrogram (STFT)}, select the chirp-plus-burst signal and step the
\emph{Window length} from 32 to 1024 samples. Watch the chirp's stripe narrow and then widen again,
while the tone burst's edges go from sharp to blurred. The sweet spot sits near $1/\sqrt\alpha$.
Then open Lab~1's \emph{A live waterfall} to see the same display running continuously.
\end{tryit}

%======================================================================
\section{Nonlinear systems in brief}
\label{sec:ch02:nonlinear}
%======================================================================

Turn a guitar amplifier up until it ``breaks up'' and a clean note becomes a gritty, buzzing one.
Nothing was added from outside: the amplifier itself manufactured new frequencies, because it was
pushed beyond the range where it behaves linearly. Guitarists pay for that sound. Radio engineers
spend their careers fighting it.

Everything above assumed linearity, and the most important consequence of linearity is that an LTI
system cannot create new frequencies. Real amplifiers, mixers and data converters are only
approximately linear, and their nonlinearity creates new frequencies---some harmless, some
devastating. Chapter~\ref{ch:sdr} treats the radio-frequency consequences in detail; here we see the
spectral mechanism.

A memoryless, weakly nonlinear device can be modelled by a power series,
\begin{equation}
y(t)=a_1x(t)+a_2x^2(t)+a_3x^3(t)+\cdots .
\label{eq:ch02:powerseries}
\end{equation}
With a single tone $x=A\cos(2\pi f_1t)$, the squared term produces DC and a second harmonic at
$2f_1$, and the cubed term produces a third harmonic at $3f_1$ and a component at $f_1$ itself,
$\tfrac34a_3A^3\cos(2\pi f_1t)$. In most amplifiers $a_3$ has the opposite sign to $a_1$, so this term
reduces the gain at large amplitudes: \term{gain compression} (Figure~\ref{fig:ch02_compression}). The
\emph{1\,dB compression point} $P_{1\text{dB}}$ is the input level at which the gain has dropped by
1\,dB.

\mdcfig{ch02_compression}{An amplifier that saturates. Left: a sine wave through a smooth limiter,
$y=\tanh(x)$, at three drive levels; hard drive flattens it towards a square wave. Right: the odd
harmonics this creates, relative to the fundamental. Doubling the drive can raise the third harmonic
by tens of decibels.}

Harmonics are far from the signal and easily filtered. The real damage comes from two or more tones
together. With $x=A\cos(2\pi f_1t)+A\cos(2\pi f_2t)$, the cubed term produces, among other things,
\term{third-order intermodulation} (IM3) products at $2f_1-f_2$ and $2f_2-f_1$
(Figure~\ref{fig:ch02_intermod}). If $f_1$ and $f_2$ are close together, these products fall right next
to them---in the passband, where no filter can remove them. A modulated signal is a continuum of
``tones'', so its intermodulation with itself spreads power into the adjacent channels: this is the
\term{spectral regrowth} visible in Figure~\ref{fig:ch02_bandwidth_defs}, and it is what ACLR limits
measure. In a receiver, two strong interferers at $f_1$ and $f_2$ can create an IM3 product exactly on
the weak wanted signal (Figure~\ref{fig:ch02_im3_receiver}).

\mdcfig{ch02_intermod}{Left: two tones at 100 and 110\,Hz through $y=x+0.05x^2-0.12x^3$. Second-order
products appear at the difference and sum frequencies and at the harmonics; the third-order products
at $2f_1-f_2$ and $2f_2-f_1$ land right beside the tones. Right: output levels against input level.
The fundamental grows at slope 1 and compresses; the IM3 product grows at slope 3. The extrapolated
lines cross at the third-order intercept point.}

The IM3 amplitude grows as $A^3$ while the fundamental grows as $A$: on a decibel plot, slope 3 versus
slope 1, so every 1\,dB increase of input raises the IM3 products by 3\,dB and their level relative to
the tones by 2\,dB. Extrapolating both straight lines, they cross at the \term{third-order intercept
point}, specified at the input (IIP3) or the output (OIP3). It is a fictitious point---the device
compresses long before reaching it---but it summarises the nonlinearity in one number, from which the
IM3 level at any input follows:
\begin{equation}
\begin{aligned}
P_{\text{IM3}}&=3P_{\text{in}}-2\,\text{IIP3}&&\text{(input-referred, dBm)},\\
P_{\text{IM3}}-P_{\text{in}}&=-2(\text{IIP3}-P_{\text{in}})&&\text{(dBc)}.
\end{aligned}
\label{eq:ch02:ip3}
\end{equation}
For the cubic model, IIP3 lies about 9.6\,dB above the 1\,dB compression point, a rule of thumb that
holds roughly for many real amplifiers.

\begin{worked}[Reading an IP3 from a two-tone test]
A low-noise amplifier is driven with two tones at $-20$\,dBm each, 1\,MHz apart. The output shows the
tones at $-5$\,dBm (15\,dB gain) and IM3 products at $-75$\,dBm, 70\,dB below the tones. By
\eqref{eq:ch02:ip3}, $\text{IIP3}=P_{\text{in}}+70/2=-20+35=+15$\,dBm, and
$\text{OIP3}=\text{IIP3}+\text{gain}=+30$\,dBm. Now suppose two interferers at $-40$\,dBm reach this
LNA, spaced so that an IM3 product falls on a wanted channel. The products are
$2(15-(-40))=110$\,dB below the interferers, at $-150$\,dBm input-referred---below the thermal noise
in a 200\,kHz channel ($-121$\,dBm), so harmless. At $-25$\,dBm each, they would be at $-105$\,dBm,
well above it (Figure~\ref{fig:ch02_im3_receiver}). Because IM3 rises 3\,dB per dB of interferer, a
receiver's tolerance to strong neighbours collapses quickly, which is why linearity and noise figure
are traded against each other in every front end (Chapter~\ref{ch:sdr}).
\end{worked}

\mdcfig[0.8]{ch02_im3_receiver}{The worked example drawn as a spectrum. A weak wanted signal at
$-100$\,dBm sits beside two strong interferers at $-25$\,dBm each. After an LNA with an IIP3 of
$+15$\,dBm, the third-order product $2f_1-f_2$ lands at $-105$\,dBm---right on top of the wanted
signal, only 5\,dB below it, and no filter after the LNA can remove it.}

\begin{tryit}[Run a two-tone test]
In Lab~35's \emph{Two-tone test and IP3}, set \emph{Input IP3} to $+15$\,dBm and \emph{Power per
tone} to $-20$\,dBm, and confirm the IM3 products sit 70\,dB below the tones. Raise the tone power one
decibel at a time and count: the tones rise 1\,dB, the IM3 products 3\,dB. Push past compression and
see how a one-point IP3 estimate goes wrong.
\end{tryit}

\mdcfig{ch02_recap}{The chapter on one page. (1) A sinusoid is the shadow of a spinning phasor.
(2) Every signal is a sum of sinusoids. (3) Short in time means wide in frequency. (4) An LTI system
smears every input with its impulse response. (5) Stripping the carrier leaves the slow complex
envelope. (6) Every FFT looks through a window, and the window's sidelobes decide what you can
see.}

Nonlinearity is not always an enemy. Mixers are deliberately nonlinear (a mixer is, at heart, a
device whose $a_2$ term multiplies two inputs); frequency multipliers exploit harmonics; envelope
detectors (Chapter~\ref{ch:analog}) use a rectifier. And the digital predistortion of
Chapter~\ref{ch:sdr} fights an amplifier's nonlinearity with a carefully designed nonlinearity of the
opposite sign.

%======================================================================
\section{Summary}
\label{sec:ch02:summary}
%======================================================================

\begin{itemize}
\item Signals are vectors: energy is squared length, correlation is angle, orthogonal signals do not
interfere, and the best approximation is a projection. Fourier analysis is the choice of sinusoids as
coordinates, and Parseval's theorem says that the choice preserves lengths.
\item Decibels turn the enormous range of signal powers into manageable numbers and turn cascaded
gains into sums; always name the reference (dBm, dBc, dBi\ldots) and use $10\log$ for power, $20\log$
for amplitude.
\item A sinusoid is the shadow of a spinning phasor; negative frequency is the wheel turning the other
way, visible only to a receiver that sees both I and Q.
\item Periodic signals have line spectra (Fourier series); aperiodic signals have continuous spectra
(Fourier transform). Jumps cause $1/f$ spectral decay and Gibbs ringing; smoothness buys faster decay.
\item The transform properties are physics: delay is linear phase, modulation is frequency shift,
filtering is multiplication, short means wide. The uncertainty principle $\sigma_t\sigma_f\ge1/(4\pi)$
makes the last one exact, with the Gaussian as the optimum.
\item Autocorrelation and energy or power spectral density are a Fourier pair; the spectrum of random
data has the shape of the pulse's energy spectrum.
\item An LTI system is its impulse response; its output is a convolution; sinusoids pass through scaled
and phase-shifted. Distortionless transmission needs flat magnitude and linear phase. Group delay is
the delay of the information; dispersion is group delay that varies with frequency. Ideal filters
are impossible (Paley--Wiener), and sharp filters ring.

\item The Hilbert transform and the analytic signal define envelope and instantaneous frequency. Every
bandpass signal is $\re\{\tilde xe^{j2\pi f_ct}\}$; the complex envelope $\tilde x=x_I+jx_Q$ carries
everything at a rate set by bandwidth. Bandpass filtering is baseband filtering; a delay becomes a
delay plus a phase rotation.
\item Bandwidth must always be defined: 3\,dB, null-to-null, noise-equivalent, 99\% occupied, $x$-dB
and masks answer different questions.
\item The DFT samples the spectrum of a windowed block. Windows trade resolution, dynamic range and
amplitude accuracy; zero padding interpolates but never resolves; the FFT computes it all in
$O(N\log N)$. Spectrograms show how spectra change, within Gabor's uncertainty limit.
\item Nonlinearity creates harmonics and intermodulation; third-order products land beside the signal
and are characterised by the intercept point.
\end{itemize}

\begin{labbox}[Signals, spectra and the complex envelope]
Two labs accompany this chapter. Lab~35, \lab{moderncomms-labs/labs/lab35\_fourier\_spectra.py}, is
devoted to its spectral tools, with nine experiments: the \emph{Fourier series builder} (harmonic
sliders, Gibbs overshoot, duty cycle and $\sigma$-smoothing); \emph{Windows, leakage and scalloping};
\emph{Zero padding is not resolution}; \emph{The uncertainty principle} (five pulse shapes on the
$\sigma_t$--$\sigma_f$ map); \emph{Convolution in slow motion} (flip, slide, multiply, integrate,
with RC, moving-average, echo and edge-detector systems); \emph{Phase delay, group delay, dispersion}
(including Chebyshev versus Bessel pulse distortion); \emph{Five bandwidths of one signal} with PA
regrowth; \emph{The spectrogram (STFT)}; and the \emph{Two-tone test and IP3}.
Lab~1, \lab{moderncomms-labs/labs/lab01\_baseband.py}, is the companion to
Section~\ref{sec:complexenvelope}: the \emph{IQ helix}; \emph{Why two ADCs: real vs IQ sampling};
\emph{Up- and down-conversion} with LO phase and frequency errors; a complete \emph{digital
down-converter} (mix, filter, decimate) that extracts a QPSK channel from an 8\,MS/s capture with the
processing gain predicted by the noise-equivalent bandwidth; the DC spur and I/Q image of a
zero-IF receiver with a blind correction; LO phase noise; ADC dynamic range and receiver sensitivity;
and \emph{A live waterfall}. The chapter's figure script, \lab{book/figscripts/ch02\_figs.py}, is a
third laboratory: every figure in the chapter is a short function whose parameters invite
modification.
\end{labbox}

\problems
\begin{enumerate}[label=\thechapter.\arabic*]
\item \emph{(Decibels)} Express in dBm: 1\,W, 25\,mW, 1\,pW and $kT\times20$\,MHz at 290\,K. A
transmitter of $+43$\,dBm feeds a cable with 3\,dB loss and an antenna of 17\,dBi gain; what is the
EIRP in dBW and in watts?
\item Show from~\eqref{eq:ch02:fs} that the Fourier coefficients of a real signal satisfy
$c_{-k}=c_k^*$. What does this imply about the spectrum at negative frequencies? Show also that an even
real signal has real coefficients and an odd real signal purely imaginary ones.
\item \emph{(Signal space)} On $[0,1]$, find the projection of $x(t)=t^2$ onto the span of $\phi_1=1$
and $\phi_2=\sqrt3(2t-1)$, and the fraction of the energy of $x$ left in the error. Verify that the error
is orthogonal to both basis functions.
\item A 25\,MHz clock with 40\% duty cycle and negligible rise time drives a long cable. Using
\eqref{eq:ch02:pulsetrain}, find the amplitudes of the first five harmonics relative to the fundamental.
Which harmonics would vanish at exactly 50\% duty cycle, and by how many decibels does the second
harmonic change when the duty cycle goes from 50\% to 48\%?
\item Prove Parseval's theorem using the convolution and conjugation properties. Use it to show that the
energy of $\sinc(t)$ is 1, and that $\int\sinc(t-m)\sinc(t-n)\,dt=\delta_{mn}$: shifted sinc pulses are
orthonormal.
\item A rectangular pulse of duration 1\,$\mu$s is transmitted. Compute its null-to-null bandwidth and
estimate its 99\% occupied bandwidth numerically. Repeat for a triangular pulse of the same energy and
base width 2\,$\mu$s, and for a raised-cosine (Hann) pulse of the same duration. Explain the ranking in
terms of smoothness.
\item \emph{(Uncertainty)} Compute $\sigma_t$, $\sigma_f$ and their product for the Gaussian
$x(t)=e^{-\pi t^2/T^2}$ and verify that the bound~\eqref{eq:ch02:uncertainty} is met with equality. Then
show that for the one-sided exponential $e^{-t}u(t)$ the RMS bandwidth is infinite, and explain why in
terms of the behaviour of $x(t)$ at $t=0$.
\item A channel has $H(f)=1+\alpha e^{-j2\pi f\tau}$ (a direct path plus one echo). Sketch
$|H(f)|$ and the group delay for $\alpha=0.5$ and $\tau=1\,\mu$s. At which frequencies are the
notches, and what is their depth? Repeat for $\alpha=1.5$; which case is minimum phase, and how can you
tell from the group delay?
\item For the $RC$ filter of the worked example in Section~\ref{sec:lti}, compute the phase delay and
group delay at $f=0$, $f_3$ and $3f_3$. Pass a 1\,MHz tone burst with a slowly varying envelope through
it and predict, from~\eqref{eq:ch02:gdpacket}, the delay of the envelope and of the zero crossings.
\item Show that the Hilbert transform of $\cos(2\pi f_0 t)$ is $\sin(2\pi f_0 t)$, and find
the analytic signal and complex envelope (with respect to $f_c$) of
$x(t)=\cos(2\pi f_0 t)$ when $f_0=f_c+\Delta$. Then find the complex envelope of
$x(t)=\cos(2\pi f_ct)\cos(2\pi f_mt)+\sin(2\pi f_ct)\sin(2\pi f_mt)$ and sketch its spectrum. Is it
symmetric?
\item A bandpass filter centred at $f_c$ has frequency response $H(f)=1/(1+j(f-f_c)/B)$ for $f>0$
(and the Hermitian extension for $f<0$). Find the baseband-equivalent impulse response $\tilde h(t)$.
Is it real? The filter is now mistuned so that its centre is at $f_c+B/2$; find the new $\tilde h(t)$ and
explain, in terms of I/Q cross-coupling, what it does to a QPSK constellation.
\item A 1024-point FFT is computed on data sampled at 1.024\,MHz with a Hann window. What
is the bin spacing? What is the noise-equivalent bandwidth of each bin? By how much does the noise
floor per bin drop if the FFT length is doubled? A tone of amplitude 1\,V lies exactly halfway between
two bins: what amplitude does the largest bin show, and how would a flat-top window change this?
\item An amplifier has $\text{IIP3}=+5$\,dBm and 20\,dB gain. Two tones at $-25$\,dBm each are applied.
Find the output tone level, the output IM3 level and the IM3 level in dBc. At what input level per tone
would the IM3 products equal the output noise floor of $-95$\,dBm?
\item \emph{(Simulation)} Using the figure script or Lab~35, find the minimum tone separation that
each of the four windows can resolve when the tones are equal in amplitude, and when one
is 40\,dB weaker. Explain the results in terms of main-lobe width and sidelobe level, and compare with
Table~\ref{tab:ch02:windows}.
\item \emph{(Simulation)} Generate a 50\,MHz, 50\%-duty trapezoidal clock with 300\,ps edges at a
simulation rate of 100\,GS/s, compute its spectrum, and compare the harmonic levels with the
two-corner envelope of Figure~\ref{fig:ch02_clock_emi}. Then add 0.5\% down-spread triangular
spread-spectrum modulation at 32\,kHz and measure the reduction of the peak level of the 11th harmonic
in a 120\,kHz resolution bandwidth.
\item \emph{(Simulation)} Write a function that computes the STFT of a signal and plots its spectrogram.
Apply it to a linear chirp of rate $\alpha=10$\,kHz/s sampled at 48\,kHz with window lengths from 1 to
100\,ms, and measure the width of the chirp's ridge. Show that the narrowest ridge occurs near
$T=1/\sqrt\alpha$.
\end{enumerate}

\furtherreading
\begin{itemize}
\item A.~V.~Oppenheim and A.~S.~Willsky, with S.~H.~Nawab, \emph{Signals and Systems}, 2nd ed.,
Prentice Hall, 1997. The standard text for continuous and discrete time side by side; thorough on
convolution, Fourier series and transforms, and sampling.
\item R.~N.~Bracewell, \emph{The Fourier Transform and Its Applications}, 3rd ed., McGraw-Hill, 2000.
Physical intuition for every transform pair, with an unequalled pictorial dictionary of transforms and
a lucid treatment of impulses, convolution and the uncertainty relation.
\item B.~P.~Lathi and R.~Green, \emph{Linear Systems and Signals}, 3rd ed., Oxford University Press,
2018. Patient and example-rich; a good first book for readers without a signals background.
\item S.~Haykin and B.~Van~Veen, \emph{Signals and Systems}, 2nd ed., Wiley, 2003; and S.~Haykin and
M.~Moher, \emph{Communication Systems}, 5th ed., Wiley, 2009. The latter's early chapters develop
Fourier analysis, the Hilbert transform and the complex envelope exactly for the communications
engineer.
\item R.~G.~Lyons, \emph{Understanding Digital Signal Processing}, 3rd ed., Prentice Hall, 2010.
Exceptionally clear on the DFT, leakage, windows, scalloping and complex (quadrature) signals; full of
practical tricks.
\item F.~J.~Harris, ``On the use of windows for harmonic analysis with the discrete Fourier
transform,'' \emph{Proceedings of the IEEE}, vol.~66, no.~1, pp.~51--83, 1978. The definitive window
reference, with the figures of merit used in Table~\ref{tab:ch02:windows}.
\item J.~W.~Cooley and J.~W.~Tukey, ``An algorithm for the machine calculation of complex Fourier
series,'' \emph{Mathematics of Computation}, vol.~19, no.~90, pp.~297--301, 1965; and M.~T.~Heideman,
D.~H.~Johnson and C.~S.~Burrus, ``Gauss and the history of the fast Fourier transform,'' \emph{IEEE ASSP
Magazine}, vol.~1, no.~4, pp.~14--21, 1984.
\item D.~Gabor, ``Theory of communication,'' \emph{Journal of the IEE, Part III}, vol.~93, no.~26,
pp.~429--457, 1946. The origin of the analytic signal and of time--frequency analysis; still readable.
\item Keysight Technologies, \emph{Spectrum Analysis Basics}, Application Note 150 (regularly updated).
The practical guide to RBW, VBW, sweep time, detectors and noise-marker corrections.
\item C.~R.~Paul, \emph{Introduction to Electromagnetic Compatibility}, 2nd ed., Wiley, 2006. The
trapezoidal-clock spectrum bounds, emission limits and the Fourier analysis behind EMC design.
\end{itemize}
